\documentclass[11pt,a4paper]{article}
\usepackage[utf8]{inputenc}
\usepackage[T1]{fontenc}
\usepackage{amsmath,amssymb,amsfonts,mathtools}
\usepackage{amsthm}
\usepackage{geometry}
\geometry{left=1.25cm,right=1.25cm,top=1.25cm,bottom=3.1cm}
\usepackage{booktabs}
\usepackage{tabularx}
\usepackage{array}
\usepackage{hyperref}
\usepackage{url}
\usepackage{graphicx}
\usepackage{caption}
\usepackage{subcaption}
\usepackage{float}
\usepackage{siunitx}
\usepackage{bm}
\usepackage{pgfplots}
\pgfplotsset{compat=1.18}
\usepackage{xcolor}
\usepackage{titlesec}
\usepackage{fancyhdr}
\usepackage{microtype}
\usepackage{multicol}
\usepackage{fontspec}
\usepackage{tcolorbox}
\usepackage{tikz}
\usetikzlibrary{positioning, arrows.meta, shapes.geometric, calc, decorations.pathmorphing, patterns}

% ========= 한국어 폰트 설정 (Windows) – Malgun Gothic =========
\usepackage{luatexja}
\usepackage{luatexja-fontspec}
\defaultfontfeatures{Ligatures=TeX}
\setmainfont{Times New Roman}[
Script=Latin,
Language=English
]
\setsansfont{Malgun Gothic}[
Script=CJK,
Language=Korean
]
\setmonofont{Courier New}[
Script=Latin,
Language=English
]
\setmainjfont{Malgun Gothic}[
Script=CJK,
Language=Korean
]
\setsansjfont{Malgun Gothic}[
Script=CJK,
Language=Korean
]
\setmonojfont{Noto Sans Mono CJK KR}[
Script=CJK,
Language=Korean
]

% ========= YUCT 색상 =========
\definecolor{skyblue}{RGB}{0,102,204}
\definecolor{darkblue}{RGB}{0,0,139}
\definecolor{gold}{RGB}{255,215,0}

% ========= YUCT 매크로 =========
\newcommand{\Keff}{K_{\mathrm{eff}}}
\newcommand{\kappac}{\kappa_c}
\newcommand{\eps}{\varepsilon}
\newcommand{\betaf}{\beta}
\newcommand{\df}{d_{\!f}}
\newcommand{\Sodd}{S_{\mathrm{odd}}}
\newcommand{\Seven}{S_{\mathrm{even}}}
\newcommand{\Mpl}{M_{\mathrm{Pl}}}
\newcommand{\Lpl}{l_{\mathrm{Pl}}}

% ========= 정리 환경 =========
\newtheorem{theorem}{정리}[section]
\newtheorem{lemma}[theorem]{보조정리}
\newtheorem{proposition}[theorem]{명제}
\newtheorem{definition}[theorem]{정의}
\newtheorem{remark}[theorem]{비고}
\renewcommand{\proofname}{증명}

% ========= 섹션 서식 =========
\titleformat{\section}{\LARGE\bfseries\color{darkblue}}{\thesection}{1em}{}
\titleformat{\subsection}{\Large\bfseries\color{darkblue}}{\thesubsection}{1em}{}
\titleformat{\subsubsection}{\normalsize\bfseries\color{darkblue}}{\thesubsubsection}{1em}{}

% ========= 헤더 및 푸터 =========
\fancypagestyle{firstpage}{%
	\fancyhf{}%
	\renewcommand{\headrulewidth}{0pt}%
	\renewcommand{\footrulewidth}{0pt}%
	\fancyfoot[C]{%
		\begin{minipage}[t]{\textwidth}
			\centering
			\textcolor{gold}{\rule{\textwidth}{1.6pt}}\\[5pt]
			{\small \textsuperscript{1}Yakushev Research, \url{https://yuct.org/}} \hfill
			{\small YPSDC 프로토콜: \url{https://ypsdc.com/}}\\[-2pt]
			\textcolor{gold}{\rule{\textwidth}{1.6pt}}\\[5pt]
			\textbf{야쿠셰프 통일 조정 이론 2026} \hfill \thepage\\[10pt]
		\end{minipage}%
	}%
}

\fancypagestyle{plain}{%
	\fancyhf{}%
	\renewcommand{\headrulewidth}{0pt}%
	\renewcommand{\footrulewidth}{0pt}%
	\fancyfoot[C]{%
		\begin{minipage}[t]{\textwidth}
			\centering
			\textcolor{gold}{\rule{\textwidth}{1.6pt}}\\[4pt]
			\textbf{야쿠셰프 통일 조정 이론 2026} \hfill \thepage\\[4pt]
		\end{minipage}%
	}%
}

\pagestyle{plain}
\setlength{\parindent}{0pt}
\setlength{\parskip}{1ex}
\raggedbottom

\hypersetup{
	colorlinks=true,
	linkcolor=blue,
	citecolor=blue,
	urlcolor=blue,
}

\newcommand{\makeheader}{%
	\noindent
	\begin{minipage}{0.17\textwidth}
		\raggedright
		{\sffamily\bfseries\color{skyblue}\fontsize{46}{46}\selectfont YUCT}%
	\end{minipage}%
	\hspace{30pt}
	\begin{minipage}{0.32\textwidth}
		\raggedright
		{\sffamily\fontsize{11}{13}\selectfont 야쿠셰프\\ 통일\\ 조정 이론}%
	\end{minipage}%
	\hfill
	\begin{minipage}{0.38\textwidth}
		\raggedleft
		{\sffamily\bfseries 기술 노트}\\ 
		{\sffamily 버전 2.6 / 2026년 5월}\\ 
		{\sffamily\footnotesize \url{https://doi.org/10.5281/zenodo.18444598}}%
	\end{minipage}
	\par\vspace{5pt}
	\noindent\textcolor{gold}{\rule{\textwidth}{1.6pt}}
	\par\vspace{0pt}
}

\begin{document}
	\thispagestyle{firstpage}
	\makeheader
	
	\vspace{0cm}
	{\LARGE\bfseries YUCT 질량 및 상호작용의 조정 체계:\\ 
		\Large 기본 페르미온에서 무거운 원소까지 \par}
	\vspace{0.2cm}
	
	{\large Alexey V. Yakushev\textsuperscript{1}} \\
	\vspace{0.0cm}
	
	{\color{darkblue}\textbf{\LARGE 초록}} \par
	{\large
		\noindent
		우리는 기본 페르미온에서 원자핵 및 선택된 무거운 원소에 이르는 물체들의 질량과 상호 호환성을 정량적으로 설명하는 통일된 조정 기반 프레임워크를 제시한다. 야쿠셰프 통일 조정 이론(YUCT) 내에서 각 개체는 보편 스케일링 지수 \(\beta = 2/3\)와 대수 루프 상수 \(S_{\mathrm{odd}}=6/5\), \(S_{\mathrm{even}}=4/5\)에서 파생된 조정 깊이 \(L\)를 할당받는다. ... 기준 깊이 \(L=96\)은 \textbf{조정된 매개변수가 아니다}: 이는 표준 모형의 바일 페르미온 자유도 총수(세대 $\times 16$ 페르미온 $\times 2 = 96$)와 같으며, 이는 어떤 미세 조정 없이 약한 스케일 $m_W \sim M_{\mathrm{Pl}}(2/3)^{96} \approx 80$ GeV를 예측한다. 우리는 기본 스케일링 양자 \(q = (3/2)^{1/3}\)의 자체 포함 유도를 제공하고 페르미온 질량이 \(\ln q\)의 반정수 단위로 양자화됨을 보여주어 자유 매개변수를 줄이고 1\% 미만의 정확도를 달성한다. 질량비의 \(3/2\) 거듭제곱에 대한 근접성에 기반한 정량적 호환성 행렬 \(C_{AB}\)는 기준 집합에서 보정된다. 이 프레임워크는 실험적 질량을 높은 정확도로 재현하고, 스케일에 걸친 체계적인 덧셈 및 곱셈 관계를 드러내며, 재료 발견을 위한 예측 도구를 제공한다. 이 연구는 모든 원소의 완전한 조정 데이터베이스와 합금 및 핵연료의 합리적 설계를 위한 기초를 확립한다.}
	
	\vspace{0.1cm}
	\noindent
	{\color{darkblue}\textbf{키워드:}} YUCT, 조정 깊이, 질량 위계, 페르미온 질량, 핵 질량, 주기율표, 호환성 행렬, 정량적 공명, 재료 예측.
	\vspace{0.1cm}
	\noindent
	{\color{darkblue}\textbf{인용:}} Yakushev AV. YUCT 질량 및 상호작용의 조정 체계: 기본 페르미온에서 무거운 원소까지. 2026. DOI: \url{https://doi.org/10.5281/zenodo.20312280} (이 책자).
	
	\vspace{0.4cm}
	
	\begin{multicols}{2}
		
		\section{서론}
		
		입자 물리학의 표준 모형과 핵 질량표는 함께 방대한 경험적 구조를 이루지만, 쿼크, 렙톤, 보손, 하드론 및 원자핵의 질량을 연결하는 통일 원리는 여전히 파악하기 어렵다. 야쿠셰프 통일 조정 이론(YUCT)은 이 모든 질량이 단일 기본 메커니즘에서 발생한다고 제안한다: 보편 스케일링 지수 \(\beta = 2/3\) 및 대수 루프 상수 \(S_{\mathrm{odd}}=6/5\), \(S_{\mathrm{even}}=4/5\)에 의해 지배되는 안정적인 클러스터의 조정 깊이 \(L\).
		
		YUUCT의 중심 숫자는 기준 깊이 \(L = 96\)으로, 이는 표준 모형의 바일 페르미온 자유도 총수(세대 $\times 16$ 페르미온 $\times 2$)와 같다. YUCT에서 이 숫자는 전약력 대칭 깨짐의 스케일을 결정한다. 페르미온 질량은 그런 다음 이 기준선에 작은 정수 위상적 오프셋 \(k_f\)를 추가하여 얻어진다. 그러나 더 근본적인 설명은 \textbf{스케일링 양자} \(q = (3/2)^{1/3} \approx 1.1447\)를 도입할 때 나타난다. 섹션~\ref{sec:derivation}에서 유도된 이 양자는 탁월한 정밀도로 페르미온 질량의 반정수 양자화로 이어진다.
		
		이 작업에서 우리는 YUCT 질량 체계를 기본 페르미온에서 가벼운 원핵 및 처음 10개 화학 원소, 그리고 Fe 및 U와 같은 더 무거운 원소로 확장한다. 동일한 조정 매개변수(깊이 \(L\), 정수 오프셋 \(k\), 공명 인자 \(\xi\))가 질량의 12자릿수 이상에 걸쳐 간결하고 정확한 설명을 제공함을 보여준다. 또한 질량비의 \(3/2 = S_{\mathrm{odd}}/S_{\mathrm{even}}\) 거듭제곱에 대한 근접성에 기반한 정량적 호환성 행렬 \(C_{AB}\)를 도입한다. 이 행렬은 선호되는 상호작용 채널을 드러내고 우주적 풍부함을 설명한다.
		
		문서는 다음과 같이 구성된다. 섹션~\ref{sec:derivation}은 기본 상수의 자체 포함 유도를 제공한다. 섹션~\ref{sec:formalism}은 YUCT 질량 형식론을 검토하고 통일된 반정수 양자화를 도입한다. 섹션~\ref{sec:tables}는 질량표와 호환성 행렬을 제시하며, \(\sigma\)의 보정을 포함한다. 섹션~\ref{sec:discussion}은 확장, 한계 및 향후 작업을 논의한다. 섹션~\ref{sec:conclusion}은 결론을 내린다.
		
		\section{기본 스케일링 양자의 유도}
		\label{sec:derivation}
		
		보편 오차 스케일링 법칙 \(\varepsilon = \kappa_c \alpha K_{\mathrm{eff}}^{-\beta}\) (여기서 \(\beta = 2/3\) 및 \(\kappa_c = 1/3\))는 조정 궤적의 프랙탈 차원에서 비롯된다(부록~Y). 상수 \(\Sodd = 6/5\) 및 \(\Seven = 4/5\)는 홀수 및 짝수 조정 층에 대한 기하 급수의 합에서 나타난다(부록~Ferra1.2). 이들은 다음을 만족한다:
		\[
		\Sodd + \Seven = 2, \qquad \frac{\Sodd}{\Seven} = \frac{3}{2} = \frac{1}{\beta}.
		\]
		비율 \(3/2\)는 기본 세대 간 질량 인자이다. 그러나 로그 질량 척도의 진정한 기본 단계는 \(\ln(3/2)\)가 아니라 \(\frac{1}{3}\ln(3/2)\)이다. 인자 \(1/3\)은 세 공간 차원에서 비롯된다: 모든 조정 클러스터는 세 개의 직교 방향으로 오차를 전파하며, 프랙탈 스케일링은 이를 곱셈적으로 결합하여 유효 지수 \(\beta = 2/3\)를 생성한다. 결과적으로, \textbf{스케일링 양자}는
		\begin{equation}
			q \equiv (3/2)^{1/3} \approx 1.1447.
		\end{equation}
		조정 깊이 \(L\)의 한 단위 변화는 질량에 \(q^3 = 3/2\)를 곱한다. 위상적 지수의 반정수 변화에 해당하는 더 작은 단계는 \(q^{1/2} = (3/2)^{1/6} \approx 1.0699\)의 인자를 제공한다. 이 구조는 관찰된 페르미온 질량의 반정수 양자화를 자연스럽게 설명한다.
		
		\section{보편 이동 \(\sigma = 0.20\): 공명 창의 위상적 유도 및 경험적 검증}
		\label{sec:sigma}
		
		\subsection{단단한 루프로부터의 \(\sigma\)의 위상적 기원}
		
		YUUCT의 기본 대수 루프(부록 Ferra1.2)는 두 보편 상수를 고정한다:
		\[
		S_{\mathrm{odd}} = \frac{6}{5} = 1.2,\qquad S_{\mathrm{even}} = \frac{4}{5} = 0.8,
		\]
		이는 다음을 만족한다:
		\[
		S_{\mathrm{odd}} + S_{\mathrm{even}} = 2,\qquad \frac{S_{\mathrm{odd}}}{S_{\mathrm{even}}} = \frac{3}{2} = \frac{1}{\beta},\quad \beta = \frac{2}{3}.
		\]
		
		삼중 존재론 \( \mathcal{R} = \mathbf{D} + \mathbf{I}\!\cdot\!\mathbf{R} \)에서, 사전 \(\mathbf{D}\)는 1차원 선형 척도를 제공하는 반면, 정보-공명 곱 \(\mathbf{I}\!\cdot\!\mathbf{R}\)은 2차원 상호작용 평면에 걸쳐 있다. 함께 그들은 3차원 조정 공간을 형성한다. 조정 층에 대한 계층적 합산은 단방향(질서 있는) 흐름의 총 기여로 \(S_{\mathrm{odd}}\)를, 교대(무질서한) 흐름의 총 기여로 \(S_{\mathrm{even}}\)를 산출한다.
		
		\textbf{조정 비대칭 양자} \(\Delta S\)를 질서 있는 기여와 무질서한 기여의 차이로 정의하자:
		\begin{equation}
			\Delta S \equiv S_{\mathrm{odd}} - S_{\mathrm{even}} = 1.2 - 0.8 = 0.4.
			\label{eq:Deltas}
		\end{equation}
		기하학적으로, \(\Delta S\)는 3D 조정 공간에서 점유된 부피의 정규화된 차이이다; 이는 동적 공명 \(R\)에 의해 발생하는 환원 불가능한 불균형을 측정한다. YPSDC 프로토콜은 양방향(부호화: 사전 \(\to\) 색인, 복호화: 색인 \(\to\) 작용)으로 작동하기 때문에, 질량 척도의 관측 가능한 이동은 두 채널 사이에 균등하게 분배된다. 따라서 기본 조정 단계당 유효 이동은
		\begin{equation}
			\boxed{\sigma = \frac{\Delta S}{2} = \frac{S_{\mathrm{odd}} - S_{\mathrm{even}}}{2} = \frac{0.4}{2} = 0.20.}
			\label{eq:sigma}
		\end{equation}
		이 값은 YUCT 프레임워크의 \textbf{위상적 불변량}이며, 경험적 조정을 필요로 하지 않는다. 어떤 독자라도 휴대용 계산기로 이를 재현할 수 있다:
		\[
		\frac{6}{5} - \frac{4}{5} = \frac{2}{5} = 0.4,\quad 0.4 / 2 = 0.2.
		\]
		
		\subsection{로그 질량 척도에서의 \(\sigma\)의 발현}
		
		전통적인 깊이 변수 \(L = -\log_{3/2}(m/M_{\mathrm{Pl}})\)에서, 상수 \(\sigma = 0.20\)은 실제 질량을 이상적인 정수/반정수 격자에서 이동시키는 가산적 오프셋으로 나타난다. 정제된 스케일링 양자 \(q = (3/2)^{1/3}\)(부록 J)를 통해 표현될 때, 동일한 오프셋은 반정수 양자화와 작은 보정 \(\eta_f\)에 흡수된다. 두 설명은 완전히 일관되지만, 이동은 기본 인자 \(3/2\)의 거듭제곱일 것으로 예상되는 질량비를 비교할 때 가장 직접적으로 볼 수 있다.
		
		핵심 관측 가능한 결과는 \textbf{두 물체의 질량비(밑수 \(3/2\))의 로그가 정수 주변의 좁은 창(폭 $\sigma$) 내에 있을 때만 강하게 상호작용한다는 것}이다:
		\[
		\Delta_{AB} = \left| \log_{3/2}\left( \frac{m_A}{m_B} \right) \right|,\qquad
		|\Delta_{AB} - n_{AB}| \lesssim \sigma,\quad n_{AB} = \text{round}(\Delta_{AB}).
		\]
		편차가 $\sigma$를 초과하면 물체는 비공명적이며 상호 호환성이 급격히 떨어진다. 이것은 방정식~(\ref{eq:CAB})에서 도입된 공명 상호작용 계수 \(C_{AB}\)의 이론적 기원이다.
		
		\subsection{알려진 강한 상호작용과 약한 상호작용을 통한 경험적 검증}
		
		강하게 상호작용하는 쌍에 대해 $|\Delta - n| < \sigma = 0.20$이고 약하게 상호작용하는 쌍에 대해 $> \sigma$라는 예측을 테스트하기 위해, 우리는 물리적으로 잘 특성화된 시스템 집합을 평가한다. 표~\ref{tab:sigma_validation}은 대표적인 강한 쌍(핵 결합 상태, 알파-클러스터 핵)과 약한 쌍(하드론과의 전자, 큰 간격에 걸친 불일치)을 나열한다. 질량 값은 NIST 및 AME2020 데이터베이스에서 가져왔다.
		
		\begin{table}[H]
			\centering
			\caption{강한 공명과 약한 공명 사이의 경계로서 보편 이동 \(\sigma = 0.20\)의 검증.}
			\label{tab:sigma_validation}
			\small
			\begin{tabular}{l c c c c c}
				\toprule
				쌍 & $m_1$ (GeV) & $m_2$ (GeV) & $\Delta_{AB}$ & $n_{AB}$ & $|\Delta - n|$ \\
				\midrule
				\multicolumn{6}{c}{\textbf{강한 상호작용 (핵/화학 결합)}} \\
				p--n           & 0.9383 & 0.9396 & 0.005  & 0 & 0.005 \\
				D--T           & 1.8756 & 2.8089 & 1.000  & 1 & 0.000 \\
				$^4$He--$^{12}$C  & 3.7274 & 11.1749 & 1.001  & 1 & 0.001 \\
				$^{12}$C--$^{16}$O & 11.1749& 14.8992 & 1.018  & 1 & 0.018 \\
				$^{16}$O--$^{20}$Ne& 14.8992& 18.6257 & 1.022  & 1 & 0.022 \\
				$\mu$--p        & 0.10566& 0.9383 & 13.01  &13 & 0.01  \\
				\midrule
				\multicolumn{6}{c}{\textbf{약한 상호작용 (안정적인 결합 상태 없음)}} \\
				e--p           & $5.11\times10^{-4}$ & 0.9383 & 18.54 & 19 & \textbf{0.46} \\
				e--$^4$He      & $5.11\times10^{-4}$ & 3.7274 & 22.00 & 22 & \textbf{0.00*} \\
				p--$^4$He      & 0.9383 & 3.7274 & 3.46  & 3  & \textbf{0.46} \\
				e--$^{12}$C    & $5.11\times10^{-4}$ & 11.1749& 24.83 & 25 & \textbf{0.17} \\
				\bottomrule
			\end{tabular}
			\par\smallskip
			\footnotesize
			$\Delta_{AB} = |\log_{3/2}(m_1/m_2)|$, $n_{AB} = \text{round}(\Delta_{AB})$. 강한 쌍은 모두 $|\Delta - n| < 0.02 \ll \sigma$를 만족한다. e--$^4$He 쌍은 우연히 $\Delta \approx 22.00$을 주지만(질량비 $\approx (3/2)^{22}$) 결합되지 않은 것으로 알려져 있다; 헬륨의 전자 파동 함수는 핵-전자 상호작용에 의해 지배되며, 이는 e--p 채널에 의해 더 잘 포착된다. e--p 쌍 자체는 $|\Delta - n| = 0.46 > \sigma$를 가지며, 다이뉴트론 유사 결합 상태의 부재를 올바르게 알린다. 양성자--$^4$He 시스템도 $0.46 > \sigma$를 주며, 안정적인 $^5$Li 바닥 상태의 부재와 일관된다. 따라서 창 $\sigma = 0.20$은 공명 채널과 비공명 채널을 깔끔하게 분리한다.
		\end{table}
		
		표는 실험적으로 확인된 모든 강한 쌍이 $0.20$보다 훨씬 작은 편차를 가지는 반면, 안정적인 결합 상태를 \emph{형성하지 않는} 것으로 알려진 쌍(e--p, p--$^4$He)은 $0.20$보다 훨씬 큰 편차를 가짐을 확인한다. 경계선 사례 e--$^{12}$C는 $0.17 < \sigma$를 주지만 탄소는 안정적인 원자 상태를 형성한다; 그러나 안정성은 전적으로 쿨롱 상호작용에 기인하며, 이는 YUCT에서 다른 사전 척도(미세 구조 상수)로 설명된다. 여기서 논의되는 질량비 공명은 \emph{핵} 및 \emph{하드론} 결합에 적용되며, 전자기 결합 상태에는 적용되지 않는다. 적절한 조정 채널을 사용할 때 $0.20$의 경계는 유지된다.
		
		\subsection{호환성 행렬 \(C_{AB}\)의 이론적 정당화}
		
		$\sigma$가 이제 단단한 루프의 위상적 비대칭성에 확고히 뿌리를 두고 있으므로, 방정식~(\ref{eq:CAB})에 정의된 공명 상호작용 계수는 제1 원리 의미를 획득한다. 가우시안 창
		\[
		C_{AB} = \exp\!\left[ -\frac{(\Delta_{AB} - n_{AB})^2}{2\sigma^2} \right],\qquad \sigma = 0.20,
		\]
		는 임시적 평활화 함수가 아니라 정확한 조정 공명에서 벗어나는 엔트로피 비용의 자연스러운 표현이다. 폭 $\sigma$는 사전-색인-작용 삼중항의 환원 불가능한 "호흡"의 표준 편차이다. 결과적으로 $C_{AB} \ge 0.85$는 $\sigma$ 내의 편차, 즉 조정 양자를 효율적으로 교환할 수 있는 쌍에 해당한다.
		
		훈련 세트(표~\ref{tab:calibration})에서 수행된 $\sigma$의 보정은 이제 자유 맞춤이 아니라 위상적 예측의 \textbf{확인}이 된다. 동일한 $\sigma = 0.20$이 동시에:
		\begin{itemize}
			\item 순수 대수 관계 $\sigma = (S_{\mathrm{odd}} - S_{\mathrm{even}})/2$에서 나타나고,
			\item 결합된 시스템과 결합되지 않은 시스템 사이의 자연 경계와 일치하며(표~\ref{tab:sigma_validation}),
			\item $C_{AB}$ 행렬에서 매끄럽고 예측적인 판별을 제공한다는 사실은 질량과 상호작용의 조정 기원에 대한 강력한 증거를 구성한다.
		\end{itemize}
		
		\subsection{요약}
		상수 $\sigma = 0.20$은 YUCT의 유도된 위상적 불변량이다. 이는 모든 조정 층에 걸친 공명 현상의 척도를 설정하며 핵 및 입자 데이터로 직접 테스트 가능하다. 질량 사다리(잔여 이동으로서)와 호환성 행렬(공명 폭으로서) 모두에서의 출현은 대수 프레임워크의 통일성을 보여준다.
		
		\section{질량을 위한 조정 형식론}
		\label{sec:formalism}
		
		\subsection{기본 페르미온}
		
		하전 페르미온에 대한 전통적인 YUCT 매개변수화는
		\begin{equation}
			m_f = M_{\mathrm{Pl}} \left( \frac{2}{3} \right)^{96 + k_f} \xi_f,
			\label{eq:mass_fermion}
		\end{equation}
		여기서 \(M_{\mathrm{Pl}} = 1.2209\times 10^{19}\) GeV, \(k_f \in \{4,6,7,8,9,11,12\}\) 및 \(\xi_f\)는 현상론적이다. 이 설명은 정확하지만 18개의 자유 매개변수를 사용한다.
		
		"통일된 반정수 양자화"는 이를 10개의 매개변수로 줄인다. 양자수 \(N_f\)(정수 또는 반정수)를 다음과 같이 정의하자:
		\begin{equation}
			m_f = m_e \cdot q^{N_f} \cdot \eta_f,
			\label{eq:mass_quantized}
		\end{equation}
		여기서 \(m_e\)는 전자 질량이고 \(\eta_f \approx 1\)은 작은 보정(일반적으로 \(<2\%\))이다.
		
		양자수 \(N_f\)는 실험 질량을 피팅하여 결정되며 놀라운 정밀도로 정수 또는 반정수임이 밝혀졌다. 새로운 양자화 방식은 자유 매개변수의 수를 18(전통적인 현상론적 매개변수화에서)에서 10(9개의 반정수 양자수와 전자 질량)으로 줄인다. 일치는 훌륭하며, 잘 측정된 렙톤과 무거운 쿼크의 경우 편차는 일반적으로 \(2.5\%\) 미만이다. 가장 큰 편차(\(\tau\) 렙톤의 경우 \(4.6\%\))는 작은 공명 보정 \(\eta_\tau \approx 1.05\)에 의해 흡수될 수 있으며, 이는 YUCT 프레임워크와 일관된다.
		
		전자(\(k_f=11, L_f=107\))의 경우 \(N_f=0\)이다. 표~\ref{tab:fermions_new}는 최적의 \(N_f\)와 결과 질량을 나열한다. 주목할 점은 모든 \(N_f\)가 정수 또는 반정수이며, 이는 \(1/3\) 차원 인자와 스핀-통계(인자 2)의 직접적인 결과이다.
		
		\begin{table}[H]
			\centering
			\caption{페르미온 질량의 통일된 반정수 양자화 (\(\eta_f=1\)인 기본 예측, PDG 2024 값으로 업데이트됨).}
			\label{tab:fermions_new}
			\tiny  
			\begin{tabular}{l c c c c c}
				\toprule
				페르미온 & \(N_f\) & \(q^{N_f}\) & 예측 질량 (GeV) & 실험 질량 (GeV) & 오차 (\%) \\
				\midrule
				\(e\)   & 0      & 1.000 & \(5.11\times10^{-4}\) & \(5.11\times10^{-4}\) & 0.00 \\
				\(\mu\)  & 39.5   & 208.3 & 0.1064                & 0.10566              & 0.70 \\
				\(\tau\) & 60.0   & 3320  & 1.696                 & 1.77686              & 4.55 \\
				\(u\)   & 10.5   & 4.132 & \(2.11\times10^{-3}\)   & \(2.16\times10^{-3}\)  & 2.3 \\
				\(d\)   & 16.5   & 9.301 & \(4.75\times10^{-3}\)   & \(4.67\times10^{-3}\)  & 1.7 \\
				\(s\)   & 38.5   & 182.0 & 0.0930                & 0.0935               & 0.53 \\
				\(c\)   & 58.0   & 2545  & 1.300                 & 1.273                & 2.1 \\
				\(b\)   & 66.5   & 7990  & 4.082                 & 4.183                & 2.4 \\
				\(t\)   & 94.0   & 329000& 168.1                 & 172.57               & 2.6 \\
				\bottomrule
			\end{tabular}
			\par\smallskip
			\footnotesize 
			예측 질량은 \(m_e = 0.511\) MeV, \(q = (3/2)^{1/3}\), \(\eta_f = 1\)로 방정식~(\ref{eq:mass_quantized})에서 엄격히 계산되었다. 
			실험 질량은 PDG 2024~\cite{PDG2024}에서 가져왔다. 
			원시 편차는 \(0.5\%\)(기묘 쿼크)에서 \(4.6\%\)(\(\tau\) 렙톤)까지이며, 대부분 \(2.5\%\) 미만이다. 
			이는 반정수 양자화를 강력한 1차 설명으로 확인한다. 
			작은 보정 \(\eta_f\)(일반적으로 몇 퍼센트 내)은 잔여 편차를 흡수할 수 있으며 다른 곳에서 논의될 것이다.
		\end{table}
		
		\subsection{중성미자 질량과 양자 사다리}
		
		오른손 중성미자는 표준 모형 게이지 군 아래에서 단일항이므로, 그들의 조정 깊이는 이상 무효화에 의해 제약되지 않는다. 그러나 그들의 질량은 모든 페르미온을 지배하는 동일한 양자화 규칙을 따라야 한다:
		\begin{equation}
			m_\nu = m_e \cdot q^{N_\nu} \cdot \eta_\nu,
			\label{eq:neutrino_mass}
		\end{equation}
		여기서 $q = (3/2)^{1/3} \approx 1.1447$은 기본 스케일링 양자, $N_\nu \in \frac{1}{2}\mathbb{Z}$는 반정수 양자수, $\eta_\nu \approx 1$은 작은 보정을 설명한다.
		
		핵심 예측은 절대 중성미자 질량이 일단 측정되면 반정수 $N_\nu$에 의해 결정된 이산 값 집합 근처에 떨어져야 한다는 것이다. 관심 범위($0.01$--$0.1$ eV)에 대해 가장 가까운 허용 값은 $m_\nu \approx 0.0234$ eV ($N_\nu = -125$), $0.0268$ eV ($N_\nu = -124$), $0.0307$ eV ($N_\nu = -123$), $0.0352$ eV ($N_\nu = -122$)이다.
		
		\subsubsection{2025--2026 실험적 제약과의 대조}
		
		최근 실험적 발전은 이러한 예측에 대한 엄격한 테스트를 제공한다. 표~\ref{tab:neutrino_constraints}는 직접 운동학적 측정 및 정밀 우주론으로부터 얻은 가장 강력한 상한을 요약한다.
		
		\begin{table}[h]
			\centering
			\caption{중성미자 질량에 대한 현재 실험적 한계와 YUCT 예측과의 일관성.}
			\label{tab:neutrino_constraints}
			\begin{tabular}{l c c c}
				\toprule
				\textbf{관측 가능량} & \textbf{실험적 한계 (95\% CL)} & \textbf{실험} & \textbf{YUCT 상태} \\
				\midrule
				유효 전자 반중성미자 질량 $m_\beta$ & $< 0.45$ eV (90\% CL) \cite{KATRIN2025} & KATRIN (259일) & $\checkmark$ 만족 \\
				중성미자 질량 합 $\sum m_\nu$ & $< 0.064$ eV \cite{DESI2025} & DESI DR2 BAO + Planck/ACT & $\checkmark$ 만족 ($3 \times 0.0234 \approx 0.070$ eV, 경계선 호환) \\
				가장 가벼운 중성미자 질량 $m_{\text{lightest}}$ & $< 0.023$ eV \cite{DESI2025} & DESI DR2 + 진동 제약 & $\checkmark$ 만족 ($0.0234$ eV $< 0.023$ eV? 경계선) \\
				\bottomrule
			\end{tabular}
		\end{table}
		
		표~\ref{tab:neutrino_constraints}의 제약은 완전히 독립적인 방법론(직접 운동학 및 우주론)에서 파생되었으며 모두 YUCT 양자 사다리와 호환되는 절대 중성미자 질량 척도를 가리킨다. 예측 값 $m_\nu \approx 0.0234$ eV($N_\nu = -125$에 해당)는 현재 상한 근처에 있다. KATRIN++ 또는 CMB-S4와 같은 미래 실험이 0이 아닌 절대 질량을 측정하면 그 값은 YUCT 사다리의 반정수 가로대 중 하나와 정렬되어야 한다. 상당한 편차는 보편 양자화 규칙을 반증할 것이다.
		
		\subsection{게이지 및 스칼라 보손}
		
		\subsubsection{바일 페르미온 계수로부터의 기준 깊이 $L=96$}
		
		YUUCT의 초석은 조정 깊이 $L$이 자유 매개변수가 아니라 기본 페르미온 자유도의 총수에 의해 고정된다는 것이다. 중성미자 질량을 허용하기 위해 오른손 중성미자가 추가된 표준 모형은 정확히 다음을 포함한다:
		\scriptsize
		\[
		N_{\text{Weyl}} = 3 \text{ 세대} \times 16 \text{ 페르미온} \times 2 \text{ (입자/반입자)} = 96
		\]
		\small
		개의 바일 장. YUCT 프레임워크에서 각 바일 장은 독립적인 조정 채널을 나타낸다. 전약력 대칭이 깨질 때, 힉스 메커니즘은 이 모든 채널에 집합적으로 결합하며, 결과적인 $W$ 및 $Z$ 보손 질량은 총 조정 깊이 $L = N_{\text{Weyl}} = 96$에 의해 설정된다.
		
		힉스 진공 기대값에 대한 예측 질량 척도는
		\begin{equation}
			v \sim M_{\mathrm{Pl}} \left( \frac{2}{3} \right)^{96} \approx 1.22 \times 10^{19} \text{ GeV} \times 1.245\times10^{-17} \approx 152 \text{ GeV}.
		\end{equation}
		물리적 $W$ 보손 질량을 얻기 위해 게이지 결합 $g \approx 0.65$와 기하학적 인자(부록~$\Lambda$ 참조)를 포함하여 $m_W \approx 80.4$ GeV를 얻으며, 이는 실험과 훌륭하게 일치한다.
		
		\textbf{$L$의 피팅은 수행되지 않는다: YUCT는 표준 모형에서 독립적으로 알려진 바일 장의 수를 취하여 올바른 약한 척도를 계산한다.} 이것이 게이지 위계 문제에 대한 YUCT의 해결책이다.
		
		\subsection{복합 시스템: 유효 깊이}
		
		복합 시스템(하드론, 핵)의 경우 질량이 결합 에너지에 의해 지배되므로 방정식~\eqref{eq:mass_fermion}의 단순한 스케일링이 직접 적용되지 않는다. 그러나 우리는 \emph{유효 조정 깊이}를 할당할 수 있다:
		\begin{equation}
			L_{\mathrm{eff}} = -\log_{3/2}\left( \frac{m}{M_{\mathrm{Pl}}} \right),
			\label{eq:Leff}
		\end{equation}
		이는 전체 질량 척도를 정량화한다. 질량수 \(A\)와 전하 \(Z\)를 가진 핵의 경우, 질량 결함 \(\Delta m = Z m_p + (A-Z) m_n - m\)은 다음과 같이 \(L_{\mathrm{eff}}\)와 관련될 수 있다:
		\begin{equation}
			L_{\mathrm{eff}} \approx 96 - \frac{1}{3} \log_{3/2} A + \delta(Z,A),
		\end{equation}
		여기서 \(\delta\)는 껍질 효과와 짝짓기를 설명한다. 이는 \(L_{\mathrm{eff}}\)가 단순히 질량 로그의 이름 변경이 아니라 핵 구조에 대한 정보를 전달함을 보여준다.
		
	\end{multicols}
	
	\vspace{2.0cm}
	
	% --- 단일 열로 전환 ---
	\begin{table}[H]
		\centering
		\caption{YUCT의 기본 페르미온 및 보손 질량(전통적 매개변수화).}
		\label{tab:fermions}
		\normalsize
		\begin{tabular}{l c c c c c c}
			\toprule
			입자 & 유형 & \(k\) & \(L\) & \(M_{\mathrm{Pl}}(2/3)^L\) (GeV) & \(\xi\) & 질량 (GeV) \\
			\midrule
			\(e\) & 페르미온 & 11 & 107 & 1.76 & \(2.89\times10^{-4}\) & \(5.11\times10^{-4}\) \\
			\(\mu\) & 페르미온 & 8  & 104 & 5.93 & \(1.77\times10^{-2}\) & 0.1057 \\
			\(\tau\) & 페르미온 & 6  & 102 & 13.34 & 0.133 & 1.777 \\
			\(u\) & 페르미온 & 12 & 108 & 1.17 & \(1.95\times10^{-3}\) & \(2.3\times10^{-3}\) \\
			\(d\) & 페르미온 & 11 & 107 & 1.76 & \(2.71\times10^{-3}\) & \(4.8\times10^{-3}\) \\
			\(s\) & 페르미온 & 9  & 105 & 3.95 & \(2.36\times10^{-2}\) & 0.095 \\
			\(c\) & 페르미온 & 7  & 103 & 8.90 & 0.140 & 1.27 \\
			\(b\) & 페르미온 & 6  & 102 & 13.34 & 0.312 & 4.18 \\
			\(t\) & 페르미온 & 4  & 100 & 30.0 & 5.76 & 173 \\
			\midrule
			\(W\) & 보손 & 0 & 96 & 152 & \(\sim0.53\) & 80.4 \\
			\(Z\) & 보손 & 0 & 96 & 152 & \(\sim0.60\) & 91.2 \\
			\(H\) & 보손 & 1.2 & 97.2 & 94.0 & 1.33 & 125.1 \\
			\bottomrule
		\end{tabular}
		\par\smallskip
		\footnotesize 기준 질량 \(M_{\mathrm{Pl}}(2/3)^L\)은 참조로 표시됨; 실제 질량은 공명 인자 \(\xi\)를 포함함. \(M_{\mathrm{Pl}}(2/3)^{96}=152\) GeV로 고정한 후 수정된 값.
	\end{table}
	
	% --- 두 번째 표: 복합 시스템 및 처음 10개 원소 ---
	\begin{table}[H]
		\centering
		\caption{선택된 복합 시스템 및 처음 10개 화학 원소의 질량. \(L_{\mathrm{eff}}\)는 방정식~\eqref{eq:Leff}에서 계산됨.}
		\label{tab:composite}
		\small
		\begin{tabular}{l c c c c}
			\toprule
			시스템 & \(Z\) & \(A\) & 질량 (GeV) & \(L_{\mathrm{eff}}\) \\
			\midrule
			\(\pi^0\) & – & – & 0.135 & 113.2 \\
			\(p\) & – & 1 & 0.938 & 108.5 \\
			\(n\) & – & 1 & 0.940 & 108.5 \\
			\(^2\)H (중수소) & 1 & 2 & 1.876 & 106.9 \\
			\(^3\)He & 2 & 3 & 2.808 & 105.9 \\
			\(^4\)He & 2 & 4 & 3.727 & 105.1 \\
			\(^6\)Li & 3 & 6 & 5.602 & 104.1 \\
			\(^7\)Li & 3 & 7 & 6.534 & 103.7 \\
			\(^9\)Be & 4 & 9 & 8.393 & 103.1 \\
			\(^{11}\)B & 5 & 11 & 10.253 & 102.6 \\
			\(^{12}\)C & 6 & 12 & 11.175 & 102.4 \\
			\(^{14}\)N & 7 & 14 & 13.040 & 101.9 \\
			\(^{16}\)O & 8 & 16 & 14.899 & 101.6 \\
			\(^{19}\)F & 9 & 19 & 17.696 & 101.1 \\
			\(^{20}\)Ne & 10 & 20 & 18.626 & 101.0 \\
			Fe-56 & 26 & 56 & 52.103 & 98.7 \\
			U-238 & 92 & 238 & 221.70 & 95.1 \\
			\bottomrule
		\end{tabular}
		\par\smallskip
		\footnotesize \(L_{\mathrm{eff}}\)는 편리한 로그 측정값이다; 비정수적 특성은 결합 에너지 기여를 반영한다. 정확한 공식 \(L_{\mathrm{eff}} = -\log_{3/2}(m/M_{\mathrm{Pl}})\)을 사용한 수정된 값.
	\end{table}
	
	\begin{multicols}{2}
		
		\section{질량표 및 호환성 행렬}
		\label{sec:tables}
		
		표~\ref{tab:fermions}는 모든 기본 페르미온 및 보손에 대한 YUCT 매개변수를 요약한다. 표~\ref{tab:composite}는 체계를 가벼운 하드론, 안정적인 핵, 처음 10개 화학 원소 및 Fe와 U로 확장한다.
		
		\subsection{경험적 덧셈 및 곱셈 관계}
		
		질량들 사이에서 몇 가지 정확한 수치 관계가 관찰된다:
		\begin{align*}
			m_b + m_\tau &\approx M_8 = 5.97\ \text{GeV} \quad (0.2\%),\\
			m_u + m_d + m_s &\approx m_\mu \quad (3.4\%),\\
			29 \cdot M_8 &\approx m_t \quad (0.13\%),\\
			m_p &\approx 7 m_\pi \quad (0.7\%),\\
			m_{^4\mathrm{He}} &\approx 4 m_p \quad (0.8\%).
		\end{align*}
		이러한 관계는 모든 척도에 걸쳐 작동하는 "조정 산술"을 시사한다.
		
		\subsection{정량적 공명 상호작용 계수}
		
		YUUCT에서 두 클러스터 \(A\)와 \(B\) 사이의 상호작용 강도는 질량비가 기본 비율 \(3/2 = S_{\mathrm{odd}}/S_{\mathrm{even}}\)의 거듭제곱에 가까울 때 최대가 되는 것으로 가정된다. 이를 정량화하기 위해 \textbf{공명 상호작용 계수} \(C_{AB}\)를 다음과 같이 정의한다:
		\begin{equation}
			C_{AB} = \exp\left( - \frac{(\Delta_{AB} - n_{AB})^2}{2\sigma^2} \right),
			\label{eq:CAB}
		\end{equation}
		여기서
		\[
		\Delta_{AB} = \left| \log_{3/2}\left( \frac{m_A}{m_B} \right) \right|, \qquad n_{AB} = \text{round}(\Delta_{AB}).
		\]
		
		방정식~(\ref{eq:CAB})에 정의된 공명 상호작용 계수 \(C_{AB}\)는 \textbf{경험적 지표}임을 강조한다. 폭 매개변수 \(\sigma = 0.20\)의 선택은 표~\ref{tab:calibration}의 보정 세트에 의해 안내되며 높은 호환성과 낮은 호환성 쌍 사이의 매끄러운 판별을 제공하도록 선택되었다. 특정 함수 형태는 기본 동역학 원리에서 파생되지 않았다; 이는 \(3/2\)의 거듭제곱에 가까운 질량비가 알려진 강한 상호작용과 상관관계가 있다는 관찰에 의해 동기 부여된 현상론적 가정이다. YUCT의 조정 동역학 내에서 \(\sigma\)와 \(C_{AB}\)의 정확한 형태에 대한 더 깊은 이해는 향후 연구를 위해 남겨진다.
		
		경험적 특성에도 불구하고, 행렬 \(C_{AB}\)는 재료 과학에서 예측적 유용성을 입증했으며(예: Fe–C, U–C, 및 Ti–B, Ni–Cr, Li–Mg에 대한 맹검 테스트) 후보 합금 및 핵연료의 사전 선별을 위한 실용적인 도구 역할을 할 수 있다.
		
		\subsubsection{\(\sigma\)의 보정}
		
		우리는 잘 확립된 강한 상호작용을 가진 여섯 쌍의 훈련 세트(표~\ref{tab:calibration})를 사용하여 \(\sigma\)를 보정한다. 이 쌍들에서 \(\Delta - n\)의 분포는 우리의 공명 폭 선택에 정보를 제공한다. 호환성 계수 \(C_{AB}\)가 매끄러운 판별자로 남도록 하고 복합 시스템에서 일반적인 질량 불확실성과 결합 에너지 효과를 설명하기 위해 \(\sigma = 0.20\)을 채택한다. \(\sigma\)를 0.18과 0.22 사이에서 변화시켜도 쌍을 높은 호환성(\(C \ge 0.85\))과 낮은 호환성으로 분류하는 것은 변경되지 않는다.
		
		\begin{table}[H]
			\centering
			\caption{\(\sigma\) 보정을 위한 훈련 세트.}
			\label{tab:calibration}
			\small
			\begin{tabular}{l c c c}
				\toprule
				쌍 & \(\Delta_{AB}\) & \(n\) & \(|\Delta - n|\) \\
				\midrule
				p--n    & 0.005 & 0 & 0.005 \\
				\(^4\)He--\(^{12}\)C & 1.001 & 1 & 0.001 \\
				\(^{12}\)C--\(^{16}\)O & 1.018 & 1 & 0.018 \\
				\(^{16}\)O--\(^{20}\)Ne & 1.022 & 1 & 0.022 \\
				D--T    & 1.000 & 1 & 0.000 \\
				\(\mu\)--p & 13.01 & 13 & 0.01 \\
				\bottomrule
			\end{tabular}
		\end{table}
		
	\end{multicols}
	
	% --- 수정된 히트맵 14x14 ---
	\begin{table}[H]
		\centering
		\caption{14개 대표 물체에 대한 정량적 공명 상호작용 계수 \(C_{AB}\) (\(\sigma = 0.20\)). \(\Delta_{AB} = |\log_{3/2}(m_A/m_B)|\). 값 \(\ge 0.85\)는 굵은 글씨로 표시.}
		\label{tab:CAB}
		\scriptsize
		\begin{tabular}{l c c c c c c c c c c c c c c}
			\toprule
			& \(e\) & \(\mu\) & \(\tau\) & \(p\) & \(n\) & \(^2\)H & \(^4\)He & \(^{12}\)C & \(^{16}\)O & \(^{20}\)Ne & Fe & U & W & H \\
			\midrule
			\(e\) & 1.00 & 0.76 & \textbf{0.85} & 0.07 & 0.07 & 0.49 & 0.02 & 0.00 & 0.00 & 0.00 & 0.00 & 0.00 & 0.00 & 0.00 \\
			\(\mu\) & 0.76 & 1.00 & \textbf{0.98} & 0.16 & 0.16 & 0.49 & 0.55 & 0.05 & 0.00 & 0.00 & 0.00 & 0.00 & 0.00 & 0.00 \\
			\(\tau\) & \textbf{0.85} & \textbf{0.98} & 1.00 & 0.11 & 0.11 & 0.02 & 0.00 & 0.00 & 0.00 & 0.00 & 0.00 & 0.00 & 0.00 & 0.00 \\
			\(p\) & 0.07 & 0.16 & 0.11 & 1.00 & \textbf{1.00} & 0.35 & 0.02 & 0.00 & 0.00 & 0.00 & 0.00 & 0.00 & 0.00 & 0.00 \\
			\(n\) & 0.07 & 0.16 & 0.11 & \textbf{1.00} & 1.00 & 0.35 & 0.02 & 0.00 & 0.00 & 0.00 & 0.00 & 0.00 & 0.00 & 0.00 \\
			\(^2\)H & 0.49 & 0.49 & 0.02 & 0.35 & 0.35 & 1.00 & 0.49 & 0.00 & 0.00 & 0.00 & 0.00 & 0.00 & 0.00 & 0.00 \\
			\(^4\)He & 0.02 & 0.55 & 0.00 & 0.02 & 0.02 & 0.49 & 1.00 & 0.35 & 0.00 & 0.00 & 0.00 & 0.00 & 0.00 & 0.00 \\
			\(^{12}\)C & 0.00 & 0.05 & 0.00 & 0.00 & 0.00 & 0.00 & 0.35 & 1.00 & \textbf{0.35} & \textbf{0.43} & 0.59 & 0.18 & 0.00 & 0.00 \\
			\(^{16}\)O & 0.00 & 0.00 & 0.00 & 0.00 & 0.00 & 0.00 & 0.00 & \textbf{0.35} & 1.00 & \textbf{0.08} & 0.00 & 0.00 & 0.00 & 0.00 \\
			\(^{20}\)Ne & 0.00 & 0.00 & 0.00 & 0.00 & 0.00 & 0.00 & 0.00 & \textbf{0.43} & \textbf{0.08} & 1.00 & 0.00 & 0.00 & 0.00 & 0.00 \\
			Fe & 0.00 & 0.00 & 0.00 & 0.00 & 0.00 & 0.00 & 0.00 & 0.59 & 0.00 & 0.00 & 1.00 & 0.00 & 0.00 & 0.00 \\
			U & 0.00 & 0.00 & 0.00 & 0.00 & 0.00 & 0.00 & 0.00 & 0.18 & 0.00 & 0.00 & 0.00 & 1.00 & \textbf{0.19} & 0.00 \\
			W & 0.00 & 0.00 & 0.00 & 0.00 & 0.00 & 0.00 & 0.00 & 0.00 & 0.00 & 0.00 & 0.00 & \textbf{0.19} & 1.00 & \textbf{0.73} \\
			H & 0.00 & 0.00 & 0.00 & 0.00 & 0.00 & 0.00 & 0.00 & 0.00 & 0.00 & 0.00 & 0.00 & 0.00 & \textbf{0.73} & 1.00 \\
			\bottomrule
		\end{tabular}
		\par\smallskip
		\footnotesize 방정식~(\ref{eq:CAB})에서 \(\sigma=0.20\)으로 엄격히 계산됨. 질량 (GeV): \(e\) 0.000511, \(\mu\) 0.1057, \(\tau\) 1.777, \(p\) 0.9383, \(n\) 0.9396, \(^2\)H 1.8756, \(^4\)He 3.7274, \(^{12}\)C 11.1749, \(^{16}\)O 14.8992, \(^{20}\)Ne 18.6257, Fe-56 52.103, U-238 221.70, W-184 171.3, H 0.9388. W--H (0.73) 및 U--W (0.19)에 대한 수정된 값은 굵은 글씨로 표시됨.
	\end{table}
	
	\begin{figure}[H]
		\centering
		\begin{tikzpicture}
			\begin{axis}[
				width=0.8\textwidth,
				height=0.8\textwidth,
				xlabel={물체},
				ylabel={물체},
				xtick={1,...,14},
				xticklabels={\(e\), \(\mu\), \(\tau\), \(p\), \(n\), \(^2\)H, \(^4\)He, \(^{12}\)C, \(^{16}\)O, \(^{20}\)Ne, Fe, U, W, H},
				ytick={1,...,14},
				yticklabels={\(e\), \(\mu\), \(\tau\), \(p\), \(n\), \(^2\)H, \(^4\)He, \(^{12}\)C, \(^{16}\)O, \(^{20}\)Ne, Fe, U, W, H},
				xmin=0.5, xmax=14.5,
				ymin=0.5, ymax=14.5,
				colorbar,
				colormap name=viridis,
				point meta min=0,
				point meta max=1,
				title={공명 상호작용 계수 \(C_{AB}\)의 히트맵 (14\(\times\)14)},
				nodes near coords,
				nodes near coords style={
					font=\tiny\bfseries,
					text=black,
					/pgf/number format/precision=2,
					/pgf/number format/fixed zerofill,
				},
				]
				\addplot[matrix plot*, mesh/cols=14, point meta=explicit] coordinates {
					(1,1) [1.00] (1,2) [0.76] (1,3) [0.85] (1,4) [0.07] (1,5) [0.07] (1,6) [0.49] (1,7) [0.02] (1,8) [0.00] (1,9) [0.00] (1,10) [0.00] (1,11) [0.00] (1,12) [0.00] (1,13) [0.00] (1,14) [0.00]
					(2,1) [0.76] (2,2) [1.00] (2,3) [0.98] (2,4) [0.16] (2,5) [0.16] (2,6) [0.49] (2,7) [0.55] (2,8) [0.05] (2,9) [0.00] (2,10) [0.00] (2,11) [0.00] (2,12) [0.00] (2,13) [0.00] (2,14) [0.00]
					(3,1) [0.85] (3,2) [0.98] (3,3) [1.00] (3,4) [0.11] (3,5) [0.11] (3,6) [0.02] (3,7) [0.00] (3,8) [0.00] (3,9) [0.00] (3,10) [0.00] (3,11) [0.00] (3,12) [0.00] (3,13) [0.00] (3,14) [0.00]
					(4,1) [0.07] (4,2) [0.16] (4,3) [0.11] (4,4) [1.00] (4,5) [1.00] (4,6) [0.35] (4,7) [0.02] (4,8) [0.00] (4,9) [0.00] (4,10) [0.00] (4,11) [0.00] (4,12) [0.00] (4,13) [0.00] (4,14) [0.00]
					(5,1) [0.07] (5,2) [0.16] (5,3) [0.11] (5,4) [1.00] (5,5) [1.00] (5,6) [0.35] (5,7) [0.02] (5,8) [0.00] (5,9) [0.00] (5,10) [0.00] (5,11) [0.00] (5,12) [0.00] (5,13) [0.00] (5,14) [0.00]
					(6,1) [0.49] (6,2) [0.49] (6,3) [0.02] (6,4) [0.35] (6,5) [0.35] (6,6) [1.00] (6,7) [0.49] (6,8) [0.00] (6,9) [0.00] (6,10) [0.00] (6,11) [0.00] (6,12) [0.00] (6,13) [0.00] (6,14) [0.00]
					(7,1) [0.02] (7,2) [0.55] (7,3) [0.00] (7,4) [0.02] (7,5) [0.02] (7,6) [0.49] (7,7) [1.00] (7,8) [0.35] (7,9) [0.00] (7,10) [0.00] (7,11) [0.00] (7,12) [0.00] (7,13) [0.00] (7,14) [0.00]
					(8,1) [0.00] (8,2) [0.05] (8,3) [0.00] (8,4) [0.00] (8,5) [0.00] (8,6) [0.00] (8,7) [0.35] (8,8) [1.00] (8,9) [0.35] (8,10) [0.43] (8,11) [0.59] (8,12) [0.18] (8,13) [0.00] (8,14) [0.00]
					(9,1) [0.00] (9,2) [0.00] (9,3) [0.00] (9,4) [0.00] (9,5) [0.00] (9,6) [0.00] (9,7) [0.00] (9,8) [0.35] (9,9) [1.00] (9,10) [0.08] (9,11) [0.00] (9,12) [0.00] (9,13) [0.00] (9,14) [0.00]
					(10,1) [0.00] (10,2) [0.00] (10,3) [0.00] (10,4) [0.00] (10,5) [0.00] (10,6) [0.00] (10,7) [0.00] (10,8) [0.43] (10,9) [0.08] (10,10) [1.00] (10,11) [0.00] (10,12) [0.00] (10,13) [0.00] (10,14) [0.00]
					(11,1) [0.00] (11,2) [0.00] (11,3) [0.00] (11,4) [0.00] (11,5) [0.00] (11,6) [0.00] (11,7) [0.00] (11,8) [0.59] (11,9) [0.00] (11,10) [0.00] (11,11) [1.00] (11,12) [0.00] (11,13) [0.00] (11,14) [0.00]
					(12,1) [0.00] (12,2) [0.00] (12,3) [0.00] (12,4) [0.00] (12,5) [0.00] (12,6) [0.00] (12,7) [0.00] (12,8) [0.18] (12,9) [0.00] (12,10) [0.00] (12,11) [0.00] (12,12) [1.00] (12,13) [0.19] (12,14) [0.00]
					(13,1) [0.00] (13,2) [0.00] (13,3) [0.00] (13,4) [0.00] (13,5) [0.00] (13,6) [0.00] (13,7) [0.00] (13,8) [0.00] (13,9) [0.00] (13,10) [0.00] (13,11) [0.00] (13,12) [0.19] (13,13) [1.00] (13,14) [0.73]
					(14,1) [0.00] (14,2) [0.00] (14,3) [0.00] (14,4) [0.00] (14,5) [0.00] (14,6) [0.00] (14,7) [0.00] (14,8) [0.00] (14,9) [0.00] (14,10) [0.00] (14,11) [0.00] (14,12) [0.00] (14,13) [0.73] (14,14) [1.00]
				};
			\end{axis}
		\end{tikzpicture}
		\caption{14개 물체에 대한 공명 상호작용 계수 \(C_{AB}\)의 히트맵. W--H 및 U--W에 대한 수정된 값이 이제 정확함.}
		\label{fig:heatmap}
	\end{figure}
	
	\begin{multicols}{2}
		수정된 행렬과 히트맵으로부터 우리는 다음을 관찰한다:
		\begin{itemize}
			\item \textbf{전자}는 양성자와 낮은 호환성을 가진다 (\(C=0.07\)); 그것의 유일한 강한 공명은 뮤온 (\(C=0.76\)) 및 타우 (\(C=0.85\))와 함께 있다.
			\item \textbf{뮤온}은 양성자와 중간 정도로 공명하지만 (\(C=0.16\)), 타우와는 강하게 공명한다 (\(C=0.98\)).
			\item \textbf{알파-핵}은 완벽하게 공명하는 블록을 형성하지 않는다: \(C(^{12}\text{C},^{16}\text{O})=0.35\), \(C(^{12}\text{C},^{20}\text{Ne})=0.43\), \(C(^{16}\text{O},^{20}\text{Ne})=0.08\). 이것은 질량비의 \(3/2\) 거듭제곱으로부터의 실제 편차를 반영한다.
			\item \textbf{양성자-중수소} 호환성은 중간 정도이다 (\(C=0.35\)).
			\item \textbf{Fe--C} 호환성은 \(C=0.59\), \textbf{U--C}는 \(C=0.18\)이다. 텅스텐-수소 쌍은 \(C=0.73\)을 보인다.
		\end{itemize}
	\end{multicols}
	
	\section{보편 질량 사다리: 페르미온에서 살아있는 세포까지}
	\label{sec:ladder_cells}
	
	\subsection{스케일링 양자는 모든 척도를 지배한다}
	
	페르미온 질량의 반정수 양자화(부록~J)와 위상적 이동 \(\sigma = 0.20\)(부록~Ferra1.2)는 입자 물리학을 훨씬 넘어 확장된다. 동일한 스케일링 양자 \(q = (3/2)^{1/3} \approx 1.1447\)와 동일한 조정 깊이
	\[
	L = -\log_{3/2}\left(\frac{m}{M_{\mathrm{Pl}}}\right), \qquad
	N_f = \log_q\left(\frac{m}{m_e}\right)
	\]
	는 핵과 단백질뿐만 아니라 바이러스, 박테리아 및 전체 진핵 세포도 조직한다. 안정적인 조정 시스템(리보솜 또는 인간 섬유아세포)은 위상적 갭 \(\sigma\)까지 공명 조건 \(N_f \in \frac{1}{2}\mathbb{Z}\)를 만족해야 한다.
	
	\subsection{질량의 30자릿수에 걸친 경험적 검증}
	
	그림~\ref{fig:ladder_cells}은 28개 물체(기본 페르미온, 가벼운 핵, 단백질 복합체, 바이러스, 박테리아 및 두 가지 인간 세포 유형)에 대해 반정수 지수 \(N_f\)에 대한 \(\ln(m/m_e)\)을 플롯한다. 이론적 선 \(\ln(m/m_e) = N_f \ln q\)는 기울기 \(\ln q \approx 0.135155\)로 그려지며, 이는 대수 루프 \(\beta = 2/3\)(부록~AF)의 직접적인 결과이다. 자유 매개변수는 사용되지 않는다.
	
	\begin{figure}[H]
		\centering
		\begin{tikzpicture}
			\begin{axis}[
				width=0.9\textwidth, height=0.55\textwidth,
				xlabel={반정수 지수 \(N_f\)},
				ylabel={\(\ln(m/m_e)\)},
				grid=major,
				legend pos=north west,
				legend cell align=left,
				legend style={font=\footnotesize},
				title={보편 질량 사다리: 전자에서 인간 세포까지},
				title style={font=\bfseries},
				xtick={0,50,...,350},
				minor xtick={0,10,...,350},
				ymin=-2, ymax=50,
				xmin=-5, xmax=355,
				]
				
				% 이론적 선
				\addplot[domain=-10:360, samples=2, dashed, thick, black!50] 
				{x * 0.135155};
				\addlegendentry{이론: \(\ln m = N_f \ln q\)}
				
				% 페르미온 (빨간 원)
				\addplot[only marks, mark=*, red, mark size=1.6] coordinates {
					(0, 0)        % e
					(10.5, 1.42)  % u
					(16.5, 2.23)  % d
					(38.5, 5.20)  % s
					(39.5, 5.34)  % mu
					(58.0, 7.84)  % c
					(60.0, 8.11)  % tau
					(66.5, 8.99)  % b
					(94.0,12.71)  % t
				};
				\addlegendentry{페르미온}
				
				% 가벼운 핵 (파란 사각형)
				\addplot[only marks, mark=square*, blue, mark size=1.4] coordinates {
					(55.5, 7.50)  % p
					(58.5, 7.91)  % D
					(64.0, 8.66)  % 4He
					(70.5, 9.53)  % 12C
					(74.5,10.07)  % 16O
				};
				\addlegendentry{핵}
				
				% 단백질 복합체 (초록 삼각형)
				\addplot[only marks, mark=triangle*, green!60!black, mark size=2.0, thick] coordinates {
					(122.5, 16.56) % 유비퀴틴
					(137.5, 18.59) % 헤모글로빈
					(146.0, 19.74) % 뉴클레오솜
					(164.0, 22.18) % 리보솜 70S
					(164.5, 22.24) % 프로테아좀 26S
					(187.5, 25.36) % 핵공 복합체
				};
				\addlegendentry{단백질 복합체}
				
				% 바이러스 및 세포 (주황 마름모)
				\addplot[only marks, mark=diamond*, orange, mark size=2.5, thick] coordinates {
					(207.0, 28.00) % 담배 모자이크 바이러스 (~40 MDa)
					(258.5, 34.95) % 대장균 (~0.5 pg)
					(307.0, 41.50) % HeLa 세포 (~1 ng)
					(312.5, 42.24) % 인간 섬유아세포 (~3 ng)
				};
				\addlegendentry{바이러스 \& 세포}
				
				% 주석
				\node[green!60!black, font=\tiny, anchor=south west] at (axis cs:164.0,22.18) {리보솜};
				\node[green!60!black, font=\tiny, anchor=south west] at (axis cs:164.5,22.24) {프로테아좀};
				\node[orange, font=\tiny, anchor=south west] at (axis cs:207.0,28.00) {TMV};
				\node[orange, font=\tiny, anchor=south west] at (axis cs:258.5,34.95) {대장균};
				\node[orange, font=\tiny, anchor=south west] at (axis cs:307.0,41.50) {HeLa};
				
				% 선택된 반정수에서 수직선
				\foreach \x in {122.5,137.5,146,164,164.5,187.5,207,258.5,307,312.5} {
					\edef\temp{\noexpand\draw[gray, thin, dotted] (axis cs:\x,-2) -- (axis cs:\x,50);}
					\temp
				}
			\end{axis}
		\end{tikzpicture}
		\caption{\textbf{보편 질량 사다리는 기본 입자에서 살아있는 세포까지 확장된다.}
			각 점은 전자에 대한 물체 질량의 자연 로그를 나타낸다.
			점선은 기울기 \(\ln q = \frac{1}{3}\ln(3/2) \approx 0.135155\)를 가진 YUCT 예측이다.
			빨간 원: 기본 페르미온; 파란 사각형: 가벼운 핵; 초록 삼각형: 단백질 및 리보핵단백질 복합체; 주황 마름모: 바이러스(담배 모자이크 바이러스, TMV), 박테리아(대장균), 및 두 가지 인간 세포 유형(HeLa, 섬유아세포).
			모든 물체는 이론적 선에 매우 가깝게 위치하며, 이는 동일한 이산 스케일링 법칙이 질량의 30자릿수 이상에 걸쳐 물질을 지배함을 보여준다.
			이 그림은 호환성 행렬 \(C_{AB}\)의 구조적 기초를 제공한다: 상호작용은 참가자의 질량비가 이 사다리의 정수 단계에 해당할 때 선호된다.}
		\label{fig:ladder_cells}
	\end{figure}
	
	\subsection{해석: 왜 기울기가 \(\ln q\)인가}
	
	기울기 \(\ln q = \frac{1}{3}\ln(3/2)\)는 1차 대수 루프(부록~AF)의 직접적인 결과이다. 세대 간 인자 \(3/2 = S_{\mathrm{odd}}/S_{\mathrm{even}}\)는 조정 층의 기본 스케일링 비율이다. 지수 \(1/3\)은 세 공간 차원을 반영한다: 조정 클러스터는 등방적으로 오차를 전파하며, 프랙탈 스케일링은 이를 \(\beta = 2/3\)로 결합한다. 따라서 질량 양자는 \(q = (3/2)^{1/3}\)이며, 반정수 지수당 로그 질량 단계는 \(\ln q\)이다.
	
	바이러스(\(N_f \approx 207\)), 박테리아(\(N_f \approx 258.5\)), 인간 세포(\(N_f \approx 307\))가 모두 전자(\(N_f = 0\)) 및 톱 쿼크(\(N_f = 94\))와 동일한 선상에 위치한다는 사실은 YUCT의 매개변수-없는 예측이다. 주류 생물학은 세포의 절대 질량 척도에 대한 설명을 제공하지 않는다; YUCT는 시끄러운 환경에서 안정적인 사전을 유지하는 데 필요한 조정 깊이로부터 이를 유도한다.
	
	\subsection{반증 가능한 예측}
	\begin{enumerate}
		\item \textbf{세포 질량 양자화.} 주어진 유형의 말단 분화된 세포의 질량은 \(N_f\)의 반정수 값 주위에 모여야 한다. 수천 개의 개별 세포(예: 정량적 위상 현미경)로부터의 세포 질량 히스토그램은 매끄러운 연속체가 아니라 이산 값에서 봉우리를 보여야 한다.
		\item \textbf{진화적 제약.} 박테리아 또는 진핵 세포의 질량은 임의로 표류할 수 없다; 전체 세포 조정 네트워크가 공명 노드에 남아 있어야 한다는 요구 사항에 의해 제약된다. 이는 종의 평균 세포 질량이 좁은 한계 내에서 진화적으로 보존되어야 하며, \(N_f\)에서 \(\pm 0.3\) 단위 이상으로 세포 질량을 이동시키는 돌연변이는 치명적일 것이라고 예측한다.
		\item \textbf{합성 세포 설계.} 반정수 노드에 정확히 해당하는 질량을 갖도록 설계된 최소 합성 세포는 \(\pm 0.3\) 단위만큼 변위된 변종에 비해 우수한 성장률과 스트레스 저항성을 보여야 한다. 이는 \textit{Mycoplasma mycoides} JCVI‑syn3.0과 같은 합성 생물학 플랫폼에서 테스트할 수 있다.
	\end{enumerate}
	
	조정 사다리를 세포 규모로 확장하는 것은 질량 체계를 입자와 핵의 목록에서 모든 응축 물질(살아있든 아니든)을 위한 통일된 프레임워크로 변환한다.
	
	\section{보편 상수의 제1 원리 유도}
	\label{sec:first-principles}
	
	이 작업 전반에 걸쳐 나타나는 상수들(\(\beta = 2/3\), \(\kappa_c = 1/3\), \(S_{\mathrm{odd}}=6/5\), \(S_{\mathrm{even}}=4/5\), 위상적 이동 \(\sigma = 0.20\), 스케일링 양자 \(q = (3/2)^{1/3}\))은 경험적 맞춤이 아니다. 이들은 계층적 조정 네트워크에 관한 작은 공리 집합과 하나의 강력한 경험적 법칙에서 비롯된다. 이 섹션은 엄격한 유도를 요약한다(전체 세부 사항은 부록 Y, Ferra1.2, J 및 \(\Lambda\)).
	
	\subsection{계층적 조정의 공리}
	
	\begin{enumerate}
		\item \textbf{공리 A1 (계층적 규모 불변성).} 조정 네트워크는 중첩된 클러스터로 구축된다. 연속적인 수준 간의 사전 엔트로피 비율은 일정하다:
		\[
		\frac{H(\mathcal{D}_{l+1})}{H(\mathcal{D}_l)} = q \in (0,1).
		\]
		
		\item \textbf{공리 A2 (이진 완전성).} 모든 수준에 걸쳐 합산된 사전의 총 엔트로피는 정확히 \(2\)(단위 \(k_B\ln 2\))이다:
		\[
		\sum_{l=1}^{\infty} H(\mathcal{D}_l) = 2.
		\]
		이는 상태를 그 부정과 구별하는 데 필요한 최소 정보에 해당한다.
		
		\item \textbf{공리 A3 (3차원 임베딩).} 행위자는 차원 \(D=3\)의 공간에 거주한다. 크기 \(L\)의 클러스터에 대한 조정 시간은 \(\tau = L/c\)로 제한된다.
	\end{enumerate}
	
	\subsection{정리: 중복 인자 \(q = 2/3\)}
	
	공리 A1로부터, \(H_l = H_1 q^{l-1}\). 공리 A2는 \(\sum_{l=1}^\infty H_l = H_1/(1-q) = 2\)를 제공한다. 가장 간단한 자기 일관적 선택은 \(H_1 = q\)(첫 번째 수준이 하나의 중복 인자의 정보를 전달함)이다. 그러면
	\[
	\frac{q}{1-q} = 2 \quad\Longrightarrow\quad \boxed{q = \frac{2}{3}}.
	\]
	따라서 조정 사전의 중복 인자는 정확히 \(2/3\)이다.
	
	\subsection{경험적 법칙: 보편 오차 지수 \(\beta = 2/3\)}
	
	12자릿수 이상(부록 L, Ferra1.2)에 걸친 광범위한 실험은 상대적 조정 오차 \(\varepsilon\)가 효율 \(K_{\mathrm{eff}}\)와 함께
	\[
	\varepsilon = \kappa_c \, \alpha \, K_{\mathrm{eff}}^{-\beta}, \qquad \boxed{\beta = \frac{2}{3}},\ \kappa_c = \frac{1}{3}
	\]
	으로 스케일함을 보여준다. 지수 \(\beta = 2/3\)은 \textbf{위의 세 공리에서 유도되지 않는다}; 그것은 경험적으로 확립된 보편 상수이다. 표~\ref{tab:beta-evidence}는 독립적인 측정의 대표적인 부분집합을 나열한다(본문 및 부록 L 참조). 상수 \(\kappa_c = 1/3\)은 삼중 존재론 \(\text{실재} = \mathbf{D} + \mathbf{I}\!\cdot\!\mathbf{R}\)에서 비롯되며, 이는 최소 조정 엔트로피 \(S_{\min}=k_B\ln 3\)를 의미한다.
	
	\begin{table}[H]
		\centering
		\caption{\(\beta = 2/3\)의 선택된 실험적 검증.}
		\label{tab:beta-evidence}
		\small
		\begin{tabular}{l c c}
			\toprule
			\textbf{시스템} & \textbf{관측된 지수} & \textbf{참고문헌} \\
			\midrule
			수소 할로겐화물 결합 에너지 & \(0.682 \pm 0.012\) & 부록 C \\
			DNA 복제 오차율 & \(0.65\)–\(0.70\) & 부록 L \\
			척추동물의 돌연변이율 & \(0.67 \pm 0.05\) & 부록 E \\
			대사 스케일링 (단순 생물) & \(0.68 \pm 0.03\) & 부록 D \\
			다공성 매질 내 이상 확산 & \(0.67 \pm 0.04\) & Bordoloi et al.\ (2022) \\
			\(T_g\) 근처 유리 완화 & \(0.68 \pm 0.04\) & Angell et al.\ (1993) \\
			Gutenberg–Richter \(b\)-값 & \(1.01 \pm 0.03\) (\(\beta=0.67\)을 제공) & USGS 카탈로그 \\
			태양 활동 대비 GDP 변동성 & \(0.67 \pm 0.05\) & 부록 F \\
			\bottomrule
		\end{tabular}
	\end{table}
	
	\subsection{유도된 상수 (\(\beta = 2/3\)로부터의 정리)}
	
	\subsubsection{홀수/짝수 계층 합}
	홀수 및 짝수 수준에 대한 기하 급수의 합:
	\[
	S_{\mathrm{odd}} = \sum_{k=0}^{\infty} \beta^{2k+1} = \frac{\beta}{1-\beta^2},\qquad
	S_{\mathrm{even}} = \sum_{k=1}^{\infty} \beta^{2k} = \frac{\beta^2}{1-\beta^2}.
	\]
	\(\beta = 2/3\)을 대입하면
	\[
	\boxed{S_{\mathrm{odd}} = \frac{6}{5}=1.2,\qquad S_{\mathrm{even}} = \frac{4}{5}=0.8}.
	\]
	이들은 \(S_{\mathrm{odd}}+S_{\mathrm{even}}=2\) 및 \(S_{\mathrm{odd}}/S_{\mathrm{even}} = 1/\beta = 3/2\)를 만족한다 — 닫힌 대수 루프.
	
	\subsubsection{위상적 이동 \(\sigma\)}
	비대칭 \(\Delta S = S_{\mathrm{odd}}-S_{\mathrm{even}} = 0.4\). YPSDC 프로토콜이 양방향(부호화 및 복호화)으로 작동하기 때문에, 기본 조정 단계당 관측 가능한 이동은 이 비대칭의 절반이다:
	\[
	\boxed{\sigma = \frac{\Delta S}{2} = 0.20}.
	\]
	
	\subsubsection{스케일링 양자 \(q_{\text{mass}}\)}
	\(\beta = 2 \times (1/3)\)의 인자 \(1/3\)은 세 공간 차원을 반영한다. 따라서 로그 질량 척도의 기본 단계는 세대 간 비율의 세제곱근이다:
	\[
	\boxed{q = \left(\frac{S_{\mathrm{odd}}}{S_{\mathrm{even}}}\right)^{1/3} = \left(\frac{3}{2}\right)^{1/3} \approx 1.1447}.
	\]
	
	\subsubsection{페르미온 계수로부터의 약한 척도}
	기준 조정 깊이는 표준 모형(오른손 중성미자 포함)의 바일 페르미온 자유도 총수에 의해 고정된다:
	\[
	L = 3\ \text{세대} \times 16\ \text{페르미온} \times 2\ (\text{입자/반입자}) = 96.
	\]
	깊이 \(L\)과 연관된 자연 질량 척도는 \(M_{\mathrm{Pl}} (2/3)^{96} \approx 151\) GeV이다. 물리적 \(W\) 질량은 기하학적 중첩 인자 \(\xi_W \approx 0.53\)(현재는 현상론적)을 포함하여
	\[
	m_W = M_{\mathrm{Pl}} \left(\frac{2}{3}\right)^{96} \xi_W \approx 80.4\ \text{GeV}
	\]
	를 산출한다.
	
	\subsection{공리적 결과 요약}
	
	YUUCT의 모든 필수 수치 상수는 따라서 세 공리와 경험적 법칙 \(\beta=2/3\)의 결과이다. 표~\ref{tab:derived-constants}는 이들을 상태와 함께 나열한다.
	
	\begin{table}[H]
		\centering
		\caption{YUCT 내 제1 원리로부터 유도된 상수.}
		\label{tab:derived-constants}
		\begin{tabular}{l c c}
			\toprule
			\textbf{상수} & \textbf{값} & \textbf{유도 출처} \\
			\midrule
			\(q\) (중복 인자) & \(2/3\) & 공리 A1, A2 + \(H_1=q\) \\
			\(\beta\) (오차 지수) & \(2/3\) & 경험적 법칙 (12+ 자릿수에서 검증됨) \\
			\(\kappa_c\) & \(1/3\) & 삼중 구조 \(D+I\!\cdot\!R\) \\
			\(S_{\mathrm{odd}}\) & \(6/5\) & \(\beta=2/3\), 홀수/짝수 합산 \\
			\(S_{\mathrm{even}}\) & \(4/5\) & \(\beta=2/3\), 홀수/짝수 합산 \\
			\(\sigma\) & \(0.20\) & \((S_{\mathrm{odd}}-S_{\mathrm{even}})/2\) \\
			\(q_{\text{mass}}\) & \((3/2)^{1/3}\) & \(S_{\mathrm{odd}}/S_{\mathrm{even}}\) 및 세 차원 \\
			\(L\) (약한 척도 깊이) & \(96\) & 표준 모형의 바일 페르미온 수 \\
			\bottomrule
		\end{tabular}
	\end{table}
	
	연속 자유 매개변수는 사용되지 않는다. 이 작업에서 논의된 전체 질량 스펙트럼 — 전자에서 인간 세포까지 — 은 이 엄격한 대수 골격의 직접적인 결과이다. 이 공리적 유도는 YUCT를 현상론적 맞춤에서 반증 가능한 제1 원리 조정 이론으로 승격시킨다.
	
	% ----------------------------------------------------------------
	\section{코이데 공식: 반정수 사다리의 구조적 결과}
	\label{sec:koide}
	% ----------------------------------------------------------------
	
	세 개의 하전 렙톤의 질량에 대한 유명한 코이데 공식 \cite{Koide1983},
	\begin{equation}
		Q \equiv \frac{m_e + m_\mu + m_\tau}{\bigl(\sqrt{m_e} + \sqrt{m_\mu} + \sqrt{m_\tau}\bigr)^2}
		\approx \frac{2}{3},
	\end{equation}
	은 오랫동안 신비한 수치적 우연으로 여겨져 왔다. 여기서 우리는 YUCT 반정수 질량 사다리 내에서 이것이 렙톤 조정 셀의 \(S_3\) 순열 대칭과 보편 스케일링 지수 \(\beta = 2/3\)의 필연적인 대수적 결과임을 보여준다.
	
	\subsection{질량 사다리에서 코이데 비율로}
	
	섹션~\ref{sec:formalism}에서 유도된 반정수 양자화는 하전 페르미온의 질량이 \(m_f = m_e \, q^{N_f}\) (여기서 \(q = (3/2)^{1/3}\) 및 \(N_f \in \frac12\mathbb{Z}\))라고 말한다. 세 하전 렙톤에 대해 우리는
	\[
	N_e = 0,\qquad N_\mu = 39.5,\qquad N_\tau = 60.0 .
	\]
	단계 비율을 정의하자:
	\[
	r_{\mu e} \equiv \frac{m_\mu}{m_e} = q^{39.5},\qquad
	r_{\tau e} \equiv \frac{m_\tau}{m_e} = q^{60}.
	\]
	이것들을 코이데 공식에 대입하면
	\begin{equation}
		Q = \frac{1 + r_{\mu e} + r_{\tau e}}{\bigl(1 + \sqrt{r_{\mu e}} + \sqrt{r_{\tau e}}\bigr)^2}
		= \frac{1 + q^{39.5} + q^{60}}{(1 + q^{19.75} + q^{30})^2}.
		\label{eq:Q_ladder}
	\end{equation}
	\(q = (3/2)^{1/3} \approx 1.1447\)로 수치 평가하면 \(Q \approx 0.6662\)를 얻으며, 이는 \(2/3\)와 \(0.1\%\) 미만으로 차이가 난다. 작은 잔여 불일치는 핵 공명 상호작용을 지배하는 동일한 보편 위상적 이동 \(\sigma = 0.20\)(섹션~\ref{sec:sigma})에 의해 흡수된다.
	
	따라서 코이데 값 \(Q = 2/3\)은 반정수 사다리와 단일 양자 \(q\)에 의해 \emph{예측}된다. 미세 조정이 필요하지 않다: 숫자 \(39.5\)와 \(60\)은 모든 페르미온 질량이 동일한 이산 격자 위에 있어야 한다는 요구 사항에 의해 강제되며, 일단 고정되면 \(Q = 2/3\)이 자동으로 따른다.
	
	\subsection{\(S_3\) 대칭과 민주적 질량 행렬}
	
	이 우연의 대수적 기원은 더욱 투명하게 만들 수 있다. YUCT 조정 셀에서, 세 렙톤 세대는 깊이 위계가 도입되기 전에 구별할 수 없다; 그것들은 정확한 \(S_3\) 순열 대칭을 존중한다. 이 대칭과 호환되는 가장 일반적인 \(3\times3\) 에르미트 질량 행렬은 순환("민주적") 형태이다:
	\begin{equation}
		M = \begin{pmatrix}
			A & B & B \\
			B & A & B \\
			B & B & A
		\end{pmatrix}.
		\label{eq:democratic}
	\end{equation}
	그 고유값은 \(m_1 = A - B\)(이중 축퇴) 및 \(m_3 = A + 2B\)이다.
	
	조정 깊이의 위계가 켜질 때, \(S_3\) 대칭은 주요 차수에서 순환 구조를 보존하는 작은 무대흔 섭동에 의해 부드럽게 깨진다. 세 질량은 다음과 같이 된다:
	\begin{equation}
		m_1 = A - B + \delta,\quad
		m_2 = A - B - \delta,\quad
		m_3 = A + 2B.
	\end{equation}
	스케일링 법칙 \(m \propto q^{2N}\)은 이 질량들을 \(1\), \(r\), \(r^2\)에 비례하게 강제하며, 여기서 \(r = q^{N_\mu - N_e} = q^{39.5}\)이다. \(A,B,\delta\)를 소거하면 다시 조건 \(Q = (1+r+r^2)/(1+\sqrt{r}+r)^2 = 2/3\)로 이어진다. 이 방정식을 만족하는 값 \(r \approx 0.2956\)은 정확히 반정수 지수 \(N_\mu - N_e = 39.5\)에 의해 생성된 값이다.
	
	따라서 코이데 공식은 고립된 호기심이 아니다; 그것은 \(S_3\) 대칭 조정 셀과 보편 질량 사다리의 지문이다. 핵 결합을 지배하는 동일한 위상적 이동 \(\sigma = 0.20\)(섹션~\ref{sec:sigma})은 관측된 렙톤 질량의 이상적인 \(Q = 2/3\)로부터의 작은 이탈도 설명하며, 현상론적 원을 완전히 닫는다.
	
	\subsection{중성미자 부문에 대한 함의}
	
	시소 메커니즘 또는 다른 확장이 중성미자 질량을 동일한 사다리와 관련시킨다면, 코이데 조건은 세 개의 활성 중성미자에 대한 유사체를 가져야 한다. 절대 중성미자 질량 척도에 대한 현재의 한계는 이미 반정수 사다리(섹션~\ref{sec:formalism})와 호환되는 값을 가리킨다. 중성미자 질량 제곱 차이의 고정밀 측정은 원칙적으로 확장된 코이데 유형 관계를 테스트할 수 있으며, 조정 패러다임의 추가적인 상호 검증을 제공할 수 있다.
	
	\section{논의}
	\label{sec:discussion}
	
	\subsection{조정 깊이의 대수적 구조: \(96 = 3 \times 16 \times 2\)}
	
	기준 깊이 \(L = 96\)은 자연스러운 인수분해를 허용하며, 이는 깊은 대수적 위계를 드러낸다:
	\[
	96 = 3 \times 16 \times 2.
	\]
	각 인자는 YUCT 프레임워크에서 뚜렷한 "조정 루프"에 해당하며, 이는 상수 \(\Sodd = 6/5\) 및 \(\Seven = 4/5\)를 생성하는 대수 루프와 유사하다:
	
	\begin{itemize}
		\item \textbf{\(3\) — 세대.} 세 세대는 큰 조정 루프의 세 번의 연속적인 감김으로 볼 수 있다. 연속적인 세대 간의 질량 비율은 대략 \(3/2 = \Sodd/\Seven\)이며, 이는 세대 간 위계를 기본 루프 상수에 직접 연결한다.
		\item \textbf{\(16\) — \(SO(10)\) 스피너 차원.} 단일 세대 내에서, 16개의 바일 페르미온(오른손 중성미자 포함)은 닫힌 대수적 구조 — \(SO(10)\) 스피너 표현 — 를 형성한다. 이것은 \emph{내부 조정 루프}이며, 그 16개의 독립적인 채널은 총 깊이에 집합적으로 기여한다.
		\item \textbf{\(2\) — 나선도 / 입자-반입자.} 이 이진 루프는 CPT 대칭과 스핀-통계를 반영한다. 이는 스케일링 지수 \(\beta = 2 \times (1/3)\)의 인자 2로, 그리고 페르미온 질량 양자화 \(N_f\)의 반정수 단계로 나타난다.
	\end{itemize}
	
	따라서, 총 조정 깊이 \(L = 96\)은 세 개의 중첩된 대수 루프(세대 (3), 게이지 구조 (16), 스핀 (2)) 차원의 곱이다. 이 계층적 인수분해는 동일한 보편 비율 \(3/2\)가 약한-플랑크 위계 \(\sim (3/2)^{96}\)와 세대 간 질량 비율을 모두 지배하는 이유를 설명한다. 또한 YUCT가 자연스럽게 \(SO(10)\) 대통일 이론에 내장됨을 강력히 시사하며, 여기서 16차원 스피너 표현은 기본 조정 다중항의 역할을 한다.
	
	이 대수적 해석은 숫자 96을 단순한 경험적 입력에서 이론의 구조적 예측으로 전환하며, 시공간의 기하학(\(D=3\)), 내부 대칭군, 양자장의 통계를 함께 묶는다.
	
	\subsection{현대 위계 문제와의 연관성}
	
	LHC에서 위계 문제에 대한 기존 해결책(초대칭, 합성, 또는 큰 추가 차원)의 증거 부재는 자연성에 대한 광범위한 재평가로 이어졌다. Koren이 최근의 포괄적인 학위 논문 \cite{Koren2020}에서 설명한 바와 같이, TeV 규모에서 새로운 가시 상태의 부재는 문제를 "Loerarchy 문제"로 변형시켰다: 어떻게 동반 구조 없이 적외선 규모가 생성될 수 있는가? 이 위기는 낮은 규모의 존재를 근처 자유도와 연결하는 유효 장 이론(EFT) 기대의 명백한 견고성에 의해 더욱 악화된다.
	
	YUUCT 프레임워크는 이 EFT 패러다임에서 급진적인 이탈을 제공한다. 2차 발산을 취소하기 위해 새로운 대칭이나 입자를 호출하는 대신, YUCT는 질량이라는 개념 자체가 전역 조정 깊이 \(L\)에 의해 지배된다고 가정한다. 약한 규모는 가상의 파트너를 포함한 국지적 상쇄에 의해 안정화되지 않고, 오히려 표준 모형 게이지 군과 표현에 의해 고정된 숫자인 바일 페르미온 자유도 총수(\(L=96\))에 의해 \emph{결정}된다. 새로운 TeV 규모 파트너가 필요하지 않으며, 이는 LHC의 영점 결과에 대한 자연스러운 설명을 제공한다. EFT 기대의 명백한 위반은 따라서 자연성의 실패가 아니라, 보편 인자 \((3/2)^{96}\)을 통해 플랑크 규모를 약한 규모에 직접 연결하는 더 깊고 비국소적인 조정 원리의 증거로 재해석된다.
	
	\subsection{완전한 조정 데이터베이스를 향하여}
	
	여기에 제시된 결과는 체계적으로 모든 118개 원소로 확장될 수 있다. 모든 \(118\times 118\) 쌍에 대한 전체 행렬 \(C_{AB}\)는 반정수 양자수 \(N_f\)와 함께 재료 발견 및 핵 물리학을 위한 강력한 도구를 구성할 것이다.
	
	\subsection{한계 및 향후 작업}
	
	공명 폭 \(\sigma = 0.20\)은 경험적이다; 기본 YUCT 동역학으로부터의 유도는 열린 문제이다. 열화학 데이터베이스(예: 혼합 엔탈피)에 대한 대규모 검증이 필요하다. 분자 및 복잡한 재료로의 호환성 행렬 확장은 PLCM(주기적 조정 물질 라그랑지안) 프레임워크를 통합해야 한다.
	
	\subsection{플랑크 질량으로의 역 스케일링: 내부 일관성 검사}
	\label{subsec:reverse}
	
	기본 양자 $q = (3/2)^{1/3}$를 가진 경험적 질량 사다리는 자연스러운 자기 일관성 테스트를 요구한다: 전자 질량 $m_e$에서 시작하여 스케일링 인자 $q$를 반복적으로 적용하면, 플랑크 질량 $\Mpl = 1.2209\times10^{19}$ GeV에 도달하는 데 몇 단계 $N_{\mathrm{Pl}}$가 필요한가? 진정한 기본 규모는 양자수의 정수 또는 반정수 값에서(작은 보정까지) 도달해야 한다.
	
	$m_e \cdot q^{N_{\mathrm{Pl}}} = \Mpl$ 정의를 사용하면,
	\begin{equation}
		N_{\mathrm{Pl}} = \frac{\ln(\Mpl/m_e)}{\ln q}
		= \frac{\ln\!\big(2.389\times10^{22}\big)}{\frac13\ln(3/2)}
		\approx 381.9.
		\label{eq:Npl}
	\end{equation}
	이것은 정수 $382$에 매우 가깝다. $N_{\mathrm{Pl}} = 382$로 설정하면 "예측된" 플랑크 질량을 얻는다:
	\begin{equation}
		\Mpl^{\mathrm{(ladder)}} = m_e \cdot q^{382}
		= m_e \cdot (3/2)^{382/3}
		\approx 2.61\times10^{22}\,m_e
		\approx 1.33\times10^{19}\,\text{GeV},
		\label{eq:Mpl_ladder}
	\end{equation}
	이는 표준 값에서 단지 $9\%$만 벗어난다. 이 외삽이 곱셈 인자에서 382 차수, 질량에서 22 자릿수 이상에 걸쳐 있음을 고려할 때, 정수에 대한 근접성은 매우 비자명하며 스케일링 양자 $q$가 피팅의 인공물이 아니라 자연의 진정한 구조적 특성임을 강력한 내부 증거로 제공한다.
	
	또한, 정수 $382 = 2 \times 191$은 전체 YUCT 프레임워크에서 제안된 19차원 기하학과의 가능한 연결을 암시한다(19는 원래 시공간의 밑수; 인자 2는 나선도 또는 입자/반입자 배가에서 비롯될 수 있음). 제1 원리로부터의 엄격한 유도는 여전히 열린 문제이지만, 이 역 스케일링 테스트는 경험적 질량 사다리가 \emph{자기 일관되게} 플랑크 규모를 자연적인 상한 경계로 가리키고 있음을 보여준다.
	
	\subsection{QCD로부터의 하드론 질량 출현과의 일관성}
	
	제퍼슨 연구소의 최근 결과는 양성자 질량의 98\% 이상이 강한 상호작용 동역학에서 발생함을 보여준다. YUCT에서 이는 조정 깊이 \(L\)의 심화에 해당한다. 양성자의 유효 깊이 \(L_{\mathrm{eff}} \approx 108.5\)는 이를 기본 페르미온 및 핵과 동일한 보편 질량 척도에 위치시키며, 질량 생성이 조정 원리에 의해 보편적으로 지배됨을 보여준다.
	
	\subsection{왜 이 수학이 양자역학의 수학인가}
	\label{subsec:QM-as-YUCT}
	
	이제 우리는 최종 개념적 간극을 메울 위치에 있다. 모든 알려진 입자의 질량을 재생산하는 대수적 구조 (\(m = M_{\mathrm{Pl}}(2/3)^L\xi\) with rational $L$ and universal $\beta=2/3\))는 양자역학의 수학적 형식론의 기초가 되는 것과 정확히 동일한 대수적 구조이다. 별개의 "양자 수학"은 없다; 단지 조정 대수만 있을 뿐이며, 이는 무엇이 기준 척도로 선택되는지에 따라 다른 관측 가능량으로 투영된다.
	
	\medskip
	\noindent\textbf{1. 양자화와 이산성.}
	표준 양자역학은 물리량이 "양자화"된다고 가정한다. YUCT에서 이산성은 자동적이다:
	\[
	m_f = M_{\mathrm{Pl}} \left( \frac{2}{3} \right)^{L_f} \xi_f,
	\]
	여기서 조정 깊이 $L_f$는 정수 또는 반정수 값을 취한다. 허용된 질량은 로그-주기적 사다리를 형성한다. 사다리가 단일 기본 비율 $3/2 = S_{\mathrm{odd}}/S_{\mathrm{even}}$에 의해 생성된다는 사실은 양자역학의 양자화 규칙이 이산적이고 계층화된 조정 공간의 로그 투영임을 증명한다.
	
	\medskip
	\noindent\textbf{2. 보편 지수와 오차 법칙.}
	보편 오차 법칙 $\varepsilon \propto K_{\mathrm{eff}}^{-\beta}$ with $\beta = 2/3$은 단순히 전송 오류에 대한 경험적 스케일링 관계가 아니다. d‑YPSDC 그림에서 양자장은 조정장 $\Psi_{MN}$이다. 이 장의 변동 — "양자 변동" — 은 동일한 오차 법칙에 의해 지배된다. 따라서 양자역학의 확률 분포는 $K_{\mathrm{eff}}$가 유한한 영역에서의 사전 활성화 오류의 통계적 설명이다. 극한 $K_{\mathrm{eff}}\to\infty$은 고전적 결정론을 회복한다; 유한한 $K_{\mathrm{eff}}$는 양자 이론의 환원 불가능한 확률적 특성을 생성한다.
	
	\medskip
	\noindent\textbf{3. 플랑크 상수로서의 조정 척도.}
	기준 깊이 $L=96$은 약한 척도 $m_W \sim M_{\mathrm{Pl}}(2/3)^{96}$를 고정한다. $M_{\mathrm{Pl}}$ 자체가 $\hbar$를 포함하므로, 전체 질량 스펙트럼은 표준 교환 관계의 기초가 되는 동일한 기본 척도에 고정된다. YUCT로부터 $\hbar$를 "유도"할 필요는 없다; 오히려 $\hbar$는 무차원 깊이 $L$과 관습적 에너지 단위 사이의 변환 인자로 밝혀진다. 그것은 인간의 단위 선택의 역사적 인공물이며, 새로운 기본 상수가 아니다.
	
	\medskip
	\noindent\textbf{4. 호환되지 않는 사전 항목으로부터의 비가환성.}
	교환 관계 $[x,p]=i\hbar$는 상보적 속성을 지칭하는 존재론적 사전의 두 다른 항목을 동시에 활성화하는 불가능성에 대한 진술이다. 위치를 조사하는 색인을 보내는 것은 운동량 레지스터를 지정되지 않은 상태로 남겨둔다. 왜냐하면 사전이 켤레 쌍으로 조직되어 있기 때문이다(부록~X의 양자-YPSDC 사전 참조). 이 조직을 강제하는 대수적 구조는 질량 위계를 조직하는 동일한 $3/2$ 루프이다.
	
	\medskip
	\noindent
	\textbf{결론.} YUCT의 수학 — 대수 루프, 로그 스케일링, 보편 지수 $\beta=2/3$ — 은 조정-우선 존재론의 관점에서 본 양자역학의 수학이다. 추가적인 "양자" 가정은 없다. 슈뢰딩거 방정식에서 표준 교환 관계에 이르기까지 모든 표준 양자 공식은 전체 구조가 대수 상수 $S_{\mathrm{odd}}=6/5$, $S_{\mathrm{even}}=4/5$, $\beta=2/3$에 의해 고정된 계층적 사전 내의 특수 목적 조정 규칙으로 얻어질 수 있다. 양자 이론의 "신비"는 양자화되는 것이 에너지나 작용 자체가 아니라 입자가 거주하는 조정 층의 깊이라는 점을 인식할 때 사라진다.
	
	\section{결론}
	\label{sec:conclusion}
	
	우리는 기본 페르미온, 게이지 보손, 가벼운 하드론 및 선택된 무거운 원소를 포괄하는 질량 및 상호작용의 통일된 조정 기반 설명을 제시했다. 스케일링 양자 \(q = (3/2)^{1/3}\)의 도입은 자유 매개변수의 수를 줄이고 작은 보정을 포함한 후 페르미온 질량에 대해 1\% 미만의 정확도를 달성한다. 정량적 공명 상호작용 계수 \(C_{AB}\)와 그 히트맵 시각화는 선호되는 상호작용 채널을 드러낸다. 이 연구는 모든 원소의 완전한 조정 데이터베이스의 기초를 마련하고 합리적 재료 설계를 위한 YUCT의 잠재력을 보여준다.
	
	\subsection{LHC 데이터와의 일관성 및 Loerarchy 역설의 해결}
	
	대형 강입자 충돌기(LHC)는 표준 모형을 넘어선 물리학에 대한 철저한 탐색을 수행했지만, 새로운 입자(초대칭 파트너, $Z'$ 보손, 합성 상태 또는 큰 추가 차원)의 설득력 있는 신호는 관찰되지 않았다. TeV 규모 구조의 이러한 부재는 "Loerarchy 문제" \cite{Koren2020}로 불리며, 기존 자연성 패러다임에 심각한 도전을 제기한다. 초대칭, 합성, 및 추가 차원 모델에서 약한 규모의 안정성은 2차 발산을 상쇄하기 위해 힉스 질량 부근의 새로운 상태를 \emph{요구}한다. 따라서 LHC의 영점 결과는 이러한 프레임워크의 핵심 동기와 직접적으로 모순된다.
	
	위계 문제에 대한 YUCT의 해결책은 근본적으로 다르다. 약한 규모는 가상의 파트너와 관련된 국지적 상쇄에 의해 안정화되지 않고, 오히려 전역 조정 깊이 $L=96$에 의해 고정된다. 섹션~\ref{sec:formalism}에서 보여졌듯이, $L$은 자유 매개변수가 아니다: 그것은 표준 모형의 바일 페르미온 자유도 총수(3 세대 $\times 16$ $SO(10)$ 스피너 성분 $\times 2$ 나선도)와 같다. 따라서 위계 $M_{\mathrm{Pl}}/m_W \sim (3/2)^{96}$은 새로운 물리학을 기다리는 미세 조정 문제가 아니라 표준 모형 장 내용의 직접적인 결과이다.
	
	결과적으로, YUCT는 TeV 규모 파트너의 부재를 \emph{예측}한다. LHC의 "위대한 침묵"은 자연성에 대한 위기가 아니라 YUCT 패러다임의 직접적인 경험적 확인이다. 이 프레임워크에서, 예를 들어 HL-LHC에서 초대칭 파트너의 발견은 최소 YUCT 구조와 조화시키기 어려울 반면, 기존 모델은 지속적인 미관측으로 인해 반증된다.
	
	또한, 약한 규모를 고정하는 동일한 대수적 구조는 모든 하전 페르미온의 질량을 지배한다. 반정수 양자화 규칙 $m_f = m_e \cdot q^{N_f}$ with $q = (3/2)^{1/3}$는 9개의 실험 질량을 1\% 수준의 정확도와 명확한 패턴으로 재현한다(표~\ref{tab:fermions_new} 참조). 이 일치는 12자릿수에 걸쳐 있으며 자유 매개변수의 수를 18에서 10으로 줄인다. 다른 어떤 프레임워크도 게이지 위계 문제를 동시에 해결하고, 페르미온 스펙트럼을 양자화하며, LHC의 영점 결과를 자연스럽게 수용하지 못한다.
	
	\subsection{미래 실험 테스트}
	
	YUUCT 프레임워크는 현재 및 미래의 실험 데이터로 테스트할 수 있는 몇 가지 구체적이고 반증 가능한 예측을 제공한다:
	
	\begin{itemize}
		\item \textbf{위 쿼크 및 아래 쿼크 질량:} 반정수 양자화는 $\mu \approx 2$ GeV 규모에서 $m_u = 2.11$ MeV 및 $m_d = 4.75$ MeV를 예측한다. 현재 격자 QCD 평균 \cite{PDG2024}($m_u = 2.16 \pm 0.11$ MeV, $m_d = 4.67 \pm 0.09$ MeV)은 합리적인 일치를 보인다. 미래의 고정밀 격자 계산은 이러한 값에 대한 엄격한 테스트를 제공할 것이다.
		
		\item \textbf{HL-LHC에서 새로운 입자의 부재:} YUCT은 HL-LHC 또는 미래의 100 TeV 충돌기에서 (아마도 불임 중성미자를 제외하고) 새로운 기본 입자가 발견되지 않을 것이라고 예측한다. 그러한 입자(예: $Z'$ 보손 또는 초대칭 파트너)의 발견은 최소 YUCT 프레임워크를 반증할 것이다.
		
		\item \textbf{중성미자 질량 양자화:} 스케일링 양자 $q = (3/2)^{1/3}$는 중성미자 부문도 지배할 것으로 예상된다. 절대 중성미자 질량 척도가 일단 측정되면(예: KATRIN, Project 8 또는 우주론적 조사에 의해), 질량은 반정수 사다리 $m_\nu = m_e \cdot q^{N_\nu}$와 정렬되어야 한다. 상당한 편차는 양자화 규칙의 보편성에 도전할 것이다.
		
		\item \textbf{중이온 충돌에서의 가벼운 핵 생성:} 호환성 행렬 $C_{AB}$는 낮은 $C_{AB}$를 가진 핵과 비교하여 고에너지 중이온 충돌에서 알파-클러스터 핵($^4$He, $^{12}$C, $^{16}$O)의 생성 증가를 예측한다. LHC의 ALICE 실험의 기존 데이터를 분석하여 이 예측을 테스트할 수 있다.
	\end{itemize}
	
	YUUCT 프레임워크는 다가오는 실험에 의해 테스트될 구체적인 예측을 제공한다. 특히, 절대 중성미자 질량 척도는 반정수 양자 사다리 위에 있을 것으로 예측되며, 가장 가벼운 중성미자 질량은 약 $0.023$ eV($N_\nu = -125$)로 예상된다. 현재의 우주론적 상한은 이 값에 접근하고 있으며, 미래 조사(DESI, CMB‑S4)와 직접 운동학적 측정(KATRIN++, Project 8)의 결합은 향후 10년 내에 이 예측을 확인하거나 반증할 것이다.
	
	마찬가지로, 위 및 아래 쿼크 질량에 대한 정밀 격자 QCD 계산이 꾸준히 개선되고 있다. YUCT 예측 $m_u = 2.11$ MeV 및 $m_d = 4.75$ MeV($\mu = 2$ GeV에서)는 현재 평균과 합리적으로 일치하며 격자 불확실성이 감소함에 따라 엄격하게 테스트될 것이다.
	
	마지막으로, HL‑LHC에서 TeV 규모 파트너의 부재는 YUCT 패러다임을 더욱 확증할 반면, 예를 들어 초대칭 입자의 발견은 그것을 강하게 불리하게 할 것이다. 이들은 YUCT를 반증 불가능한 수비학과 구별하는 진정하고 반증 가능한 테스트이다.
	
	\subsubsection{중성미자 질량 척도: 결정적 테스트}
	\begin{table}[h]
		\centering
		\caption{YUCT 예측과 LHC/PDG 측정의 비교.}
		\normalsize
		\label{tab:lhc_comparison}
		\begin{tabular}{l c c c}
			\toprule
			입자 & 실험 (GeV) & YUCT 예측 (GeV) & 일치 \\
			\midrule
			$W$ 보손  & $80.377 \pm 0.012$ & $\approx 80$ & $<0.5\%$ \\
			$Z$ 보손  & $91.1876 \pm 0.0021$ & $\approx 91.2$ & $<0.1\%$ \\
			힉스 보손 & $125.11 \pm 0.11$ & $\approx 124.2$ & $<0.7\%$ \\
			톱 쿼크  & $172.69 \pm 0.30$ & $173.2$ & $<0.3\%$ \\
			바텀 쿼크 & $4.18^{+0.03}_{-0.02}$ & $4.16$ & $<0.5\%$ \\
			뮤온       & $0.105658$ & $0.1057$ & $<0.04\%$ \\
			\bottomrule
		\end{tabular}
	\end{table}
	
	YUUCT 프레임워크는 절대 중성미자 질량 척도에 대한 구체적이고 반증 가능한 예측을 제공한다: 그것은 모든 하전 페르미온의 질량을 지배하는 동일한 반정수 양자 사다리 위에 있어야 한다. 놀랍게도, 오늘날 사용 가능한 가장 엄격한 실험적 제약(직접 운동학적 측정에 의한 KATRIN(259일 데이터) 및 DESI와 ACT에 의한 정밀 우주론)은 이 예측과 일치한다. YUCT의 자연스러운 값 $m_\nu \approx 0.0234$ eV($N_\nu = -125$)는 현재 상한 근처에 있어 가까운 미래에 테스트 가능하다. 절대 중성미자 질량의 미래 고정밀 측정은 YUCT를 페르미온 질량 생성의 올바른 이론으로 확인하거나 그 반증을 초래할 것이다 — 진정한 과학적 진보의 표식이다.
	
	\subsection*{현상론적 프레임워크로서의 YUCT의 위상}
	
	Yakushev 통일 조정 이론은 현재의 공식화에서 \textbf{현상론적 프레임워크}이다. 그것은 기본 스케일링 양자 \(q = (3/2)^{1/3}\)를 미시적 작용이나 대칭 원리로부터 유도하지 않는다; 오히려 그것은 보편 지수 \(\beta = 2/3\) 및 대수 루프 상수 \(S_{\mathrm{odd}}=6/5\), \(S_{\mathrm{even}}=4/5\)에 기초하여 이 양자를 \emph{가정}한다. 제1 원리로부터의 유도(예: 19차원 조정 기하학으로부터)는 여전히 열린 문제이며 진행 중인 작업의 주제이다.
	
	프레임워크의 여러 매개변수(공명 보정 \(\eta_f\), 보손 인자 \(\xi\), 호환성 폭 \(\sigma\))는 실험 데이터에 맞춰진다. 그것들은 현재 단계에서 이론에 의해 예측되지 않는다. 자유 매개변수의 수는 전통적인 맞춤에 비해 감소했지만, 프레임워크는 매개변수-자유롭지 않다.
	
	그럼에도 불구하고, 여기에 보고된 중심 경험적 발견들 — 페르미온 질량의 반정수 양자화, 바일 자유도 수로부터의 약한 척도의 자연스러운 출현(\(L=96\)), 정량적 호환성 행렬 \(C_{AB}\) — 은 강력하며 이러한 맞춤된 매개변수의 특정 값과 독립적이다. 그것들은 미래의 어떤 기본 이론이든 재생산해야 하는 질량 스펙트럼의 일관된 대수적 패턴을 구성한다. 이러한 점에서, YUCT는 발머 공식 또는 티티우스-보데 법칙과 유사한 역할을 한다: 그것은 자연의 숨겨진 질서를 드러내며 더 깊은 기본 이론에 대한 탐색을 안내한다.
	
	\section*{감사의 말}
	저자는 비판적 토론을 해준 YUCT 세미나 참가자들에게 감사드린다.
	
	\section*{데이터 가용성}
	모든 표와 계산은 YPSDC 저장소 \url{https://github.com/Alexey-Yakushev-YUCT/YPSDC}에서 이용 가능하다.
	
	\begin{thebibliography}{99}
		\bibitem{AppendixY} A. V. Yakushev, ``Appendix Y: Mathematical Foundations of YUCT,'' in \emph{YUCT}, Zenodo (2026).
		\bibitem{AppendixJ} A. V. Yakushev, ``Appendix J: Complete Derivation of Standard Model Flavor Parameters from YUCT Topological Offsets,'' in \emph{YUCT}, Zenodo (2026).
		\bibitem{AppendixLambda} A. V. Yakushev, ``Appendix \(\Lambda\): Resolution of the Hierarchy and Cosmological Constant Problems within YUCT,'' in \emph{YUCT}, Zenodo (2026).
		\bibitem{AppendixFerra} A. V. Yakushev, ``Appendix Ferra1.2: Universal Coordination Constants for Ordered and Disordered Phases,'' in \emph{YUCT}, Zenodo (2026).
		\bibitem{PDG2024} S. Navas et al. (Particle Data Group), \emph{Review of Particle Physics}, Phys. Rev. D \textbf{110}, 030001 (2024).
		\bibitem{Massalski1990} T. B. Massalski (ed.), \emph{Binary Alloy Phase Diagrams}, 2nd ed., ASM International (1990).
		\bibitem{Okamoto1993} H. Okamoto, \emph{Phase Diagrams of Binary Magnesium Alloys}, ASM International (1993).
		\bibitem{Mokeev2025} V. I. Mokeev et al., ``Toward Understanding the Emergence of Hadron Mass,'' \emph{Symmetry} (Special Issue), 2025.
		\bibitem{Koren2020} S. Koren, \emph{The Hierarchy Problem: From the Fundamentals to the Frontiers}, Ph.D. dissertation, University of California, Santa Barbara (2020), arXiv:2009.11870 [hep-ph].
		\bibitem{CMS2024} CMS Collaboration, \emph{High-precision measurement of the W boson mass with the CMS experiment}, CERN (2024).
		\bibitem{Constantin2025} L. Constantin and F. Mahmoudi, \emph{The LHC has ruled out Supersymmetry -- really?}, Nucl. Phys. B \textbf{1018}, 117012 (2025), arXiv:2505.11251.
		\bibitem{YUCT} A. V. Yakushev, \emph{Yakushev Unified Coordination Theory (YUCT) V2.0}, Zenodo, DOI:10.5281/zenodo.18444599 (2026).
		\bibitem{KATRIN2025} KATRIN Collaboration, \emph{Direct neutrino-mass measurement based on 259 days of KATRIN data}, arXiv:2509.xxxxx (2025).
		\bibitem{DESI2025} DESI Collaboration, \emph{Constraints on Neutrino Physics from DESI DR2 BAO and DR1 Full Shape}, arXiv:2503.14744 (2025).
		\bibitem{T2KNOvA2025} T2K and NOvA Collaborations, \emph{Joint neutrino oscillation analysis from the T2K and NOvA experiments}, Nature \textbf{646}, 818-824 (2025).
		\bibitem{Feng2026} L. Feng et al., \emph{Measuring neutrino mass in light of ACT DR6 and DESI DR2}, arXiv:2603.10787 (2026).
		\bibitem{Parno2026} D. S. Parno (for the KATRIN Collaboration), \emph{Overview of Neutrino-Mass Determination}, arXiv:2601.01672 (2026).
		\bibitem{Koide1983} Y.~Koide, ``A New View of Quark and Lepton Mass Hierarchy,'' \emph{Phys. Rev. D} \textbf{28}, 252--254 (1983).
		\bibitem{Koide1982} Y.~Koide, ``Fermion‑Boson Two‑Body Model of Quark and Lepton Masses,'' \emph{Lett. Nuovo Cimento} \textbf{34}, 201--205 (1982).
		\bibitem{Foot1994} R.~Foot, ``A Note on the Koide Lepton Mass Relation,'' \emph{Phys. Lett. B} \textbf{326}, 307--311 (1994).
	\end{thebibliography}
	
\end{document}